% \begin{savequote}
% {\tiny Ein Mann ging von Jerusalem nach Jericho hinab und wurde von
% Räubern überfallen. Sie plünderten ihn aus und schlugen ihn nieder; 
% dann gingen sie weg und ließen ihn 
% halbtot liegen. Zufällig kam ein Priester denselben Weg
% herab; er sah ihn und ging weiter. Auch ein Levit kam zu
% der  Stelle; er sah ihn und ging weiter. Dann kam ein Mann
% aus  Samarien, der auf der Reise war. Als er ihn sah, hatte
% er  Mitleid, ging zu ihm hin, goss Öl und Wein auf seine
% Wunden  und verband sie. 
% 
% Dann hob er ihn auf sein Reittier,
% brachte  ihn zu einer Herberge und sorgte für ihn. Am
% anderen Morgen  holte er zwei Denare hervor, gab sie dem
% Wirt und sagte:  Sorge für ihn, und wenn du mehr für ihn
% brauchst, werde  ich es dir bezahlen, wenn ich wiederkomme.
% Was meinst du:  Wer von diesen dreien hat sich als der
% Nächste dessen erwiesen,  der von den Räubern überfallen
% wurde?  
% 
% Der Gesetzeslehrer antwortete: Der, der
% barmherzig\index{Barmherzigkeit} an  ihm gehandelt hat. Da
% sagte Jesus zu ihm: Dann geh und handle genauso!
% 
% Lk 10,25-37}
% % \qauthor{{\tiny Lk 10,25-37}}
% \end{savequote}
% 
\def\ChapterCode{HKL}
\chapter
[\CO{[HKL]} Herz-Kreis\-lauf\-stör\-ung\-en]
{Herz-Kreis\-lauf\-stör\-ung\-en}
\label{chp:HKL}
\ChapterIntro
{%
Herz-Kreislauferkrankungen sind die häufigste
Todesursache.\footnote{Statistik Austria, Pressemitteilung
9.691-133/10; Statistisches Bundesamt Deutschland,
Pressemitteilung Nr.344 vom 15.09.2009} Fast jeder zweite
Todesfall geht auf eine Erkrankung des
Herz-Kreislauf-Systems zurück.  Mit steigendem Lebensalter
erhöht sich die Wahrscheinlichkeit an einer solchen zu
erkranken. 
Eine wichtige Gruppe stellen die erworbenen, durch den
Lebensstil bedingten Erkrankungen dar. Zu den wesentlichen
Risikofaktoren zählen Übergewicht, mangelnde Bewegung,
Rauchen, Bluthochdruck, erhöhte Blutfettspiegel und ein
erhöhter Blutzuckerspiegel.

Akute Notfälle ereignen sich oft nicht unvorhersehbar, oft
leidet der Patient vorab an \E{chronischen Varianten} einer
Herz-Kreislauferkrankung bzw. einer verwandten
Grunderkrankung, welche zu einem Notfallgeschehen führen
kann.
}
{%
\begin{description}
  \item[Maintainer:] Sebastian Gabriel
  \item[Autoren:] Diverse
  \item[Reviewer:] Standard-Reviewprozess
  \item[Version:] {\VarDocVersionTypeName} ({\VarDocVersionTypeDescription})
  \item[SHA1:] {\DocGitCommit}
\end{description}

{\TxTeilkapitelAASS}
}


\nocite{Harrisons18De,Harrisons18DeKap224,Harrisons18DeKap225,Harrisons18DeKap226}


\Experimental{
\Layout{57}{}{}{}{%
\begin{tabularx}{\linewidth}{XXXX}\FontTableBase
&
\TabHeading{Chronische Variante bzw. Grunderkrankung}
&
\TabHeading{Akuter Notfall}
&
\TabHeading{Anm.}
\\\addlinespace\midrule\addlinespace

&
KHK, Angina pectoris
&
Akutes Koronarsyndrom (instabile Angina pectoris,
Herzinfarkt) 
&

\\\addlinespace

&
Hypertension
&
Hypertensive Krise, Hypertensiver Notfall
&

\\\addlinespace
&
pAVK
&
Artrieller Gefäßverschluss
&

\\\addlinespace

&
Vorhofflimmern, Störungen des Reizleitungssystems, \ldots
&
Tachykarde Attacke
&
\\\addlinespace

&
Herzinsuffizienz
&
Dekompensierte Herzinsuffizienz, kard. Lungenödem
&


% \\\addlinespace
% Grad
% &
% 
% &
% 
% &

\\\addlinespace

\end{tabularx}
}
{Beispiele für den Zusammenhang von chronischen und akuten
Herz-Kreislauf-Erkrankungen}
{}

} % STATUS-EXPERIMENTAL





\label{sec:herz-kreislauf-stoerungen}
\label{topic:herz-kreislauf-stoerungen}


\clearpage
% \section{Herz-Kreislauf-Stillstand}
% 
% Der Herz-Kreislauf-Stillstand wird
% unter
% \FullRef{topic:kreislaufstillstand},
% \FullRef{topic:reanimation-eh}, und
% \prettyref{chp:BLS} und \prettyref{chp:BLS} besprochen.






%%%%%%%%%%%%%%%%%%%%%%%%%%%%%%%%%%%%%%%%%%%%%%%%%%%%%%%%%%%
%%%%%%%%%%%%%%%%%%%%%%%%%%%%%%%%%%%%%%%%%%%%%%%%%%%%%%%%%%%
%%%%%%%%%%%%%%%%%%%%%%%%%%%%%%%%%%%%%%%%%%%%%%%%%%%%%%%%%%%
%%%%%%%%%%%%%%%%%%%%%%%%%%%%%%%%%%%%%%%%%%%%%%%%%%%%%%%%%%%
\section{Störungen des \HeadingKeyword{Kreislaufs}}
\subsection{Der \HeadingKeyword{Schock} ist eine lebensbedrohliche Kreislaufstörung}
\index{Schock}
\label{topic:schock}

\TextSlide
{{\TxBeschreibung}}
{%
\DefinitionInPara{Ein Schock ist eine 
\EE{lebensbedrohliche Kreislaufstörung} aufgrund von
Unterversorgung des Körpers mit Blut und Sauerstoff infolge des
Versagens des Kreislaufsystems.} 
Es besteht eine
Missverhältnis zwischen dem Herzauswurfs und dem Bedarf des
Körpers. Dies kann unterschiedliche Ursachen haben, z.\,B.
dass das  zu wenig Herz pumpt, zu wenig Blut bzw.
Flüssigkeit vorhanden ist oder die Gefäße zu weit gestellt
sind und Blut dadurch versackt.
Es können dabei einzelne oder mehrere Teile des
Kreislaufsystems betroffen sein  (Blut, Gefäße, Herz). 

\E{Kein} Schock besteht bei kurzzeitigen
Blutverteilungsstörungen z.\,B. im Rahmen eines Kollaps
bzw. einer Synkope, wobei es auch zu Bewusstseinsstörungen
durch kurzfristige Minderdurchblutung des Gehirnskommen
kann.}
{%
\begin{itemize}
  \item Lebensbedrohliche Kreislaufstörung
  \item Kreislaufinsuffizienz 
  \item \EE{Missverhältnis
  Herz\-minuten\-volumen\,:\,Sauer\-stoff\-be\-darf}
\end{itemize}
}
{}
{}
{}


\begin{Synopsis}
\SynopsisItem{Schock = Kreislaufinsuffizienz =
  \EE{Missverhältnis
  Herz\-minuten\-volumen\,:\,Sauer\-stoff\-be\-darf} des Körpers}
\SynopsisItem{Es können einzelne oder mehrere Teile des
Kreislaufsystems betroffen sein  (Blut, Gefäße, Herz).}
\SynopsisItem{Beim Schock kommt zu wenig sauerstoffreiches Blut 
zu den Organen.}
\end{Synopsis}







\TextSlide{Kompensation, Dekompensation}
{%
\index{Schock!Kompensation}%
\index{Schock!Dekompensation}%
%%%
Der Körper versucht durch Aktivierung des Sympathikus (\FullPrettyRef{topic:sympathikus}: \Speech{\enquote{Fight!}})
gegenzusteuern (zu \E{kompensieren}). Der wichtigste
Mechanismus ist die
\T{Zentralisation}. Darunter versteht
man die \E{Blutumverteilung zugunsten lebenswichtiger Organe}
wie Herz, Gehirn, Lunge und Leber. Dabei werden z.\,B. Haut,
Muskel und Darm sowie die Extremitäten (Peripherie) kaum
mehr  durchblutet.
Periphere Venen kollabieren und die \I[Rekapillarisierungszeit!Zentralisation]{Rekapillarisierungszeit}
verlängert sich. Die Hautdurchblutung wird verringert und
die Haut wird kühl und bleich.

Dies funktioniert allerdings nur eine
gewisse Zeit und in einem gewissen Ausmaß gut, daher ist
eine schnell  einsetzende Therapie wichtig! 
Wenn die Kompensationsmaßnahmen nicht mehr ausreichen, spricht man vom \EE{{dekompensierten
Schock}}\index{Schock!dekompensierter}. Es besteht
spätestens dann akute \EE{Lebensgefahr}!
}{
\begin{itemize}
    \item {Gegensteuerung} (\EE{Kompensation})
    \item Sympathikus (\Speech{\enquote{Fight!}})
    \item Zentralisation
    \item Ende: \EE{Dekompensierter Schock}
\end{itemize}
}{}{}{}




% \Cave{Bei jedem Patienten ist zu erheben, ob ein Schock
% oder eine realistische Schockgefahr besteht oder nicht bzw. festzustellen um welche Art von
% Schock es sich handelt (s.\,u.)!}

\begin{Danger}
  \item Ein Schockzustand ist grundsätzlich als vitale
  Bedrohung einzustufen!
\end{Danger}



%%%%%%%%%%%%%%%%%%%%%%%%%%%%%%%%%%%%%%%%%%%%%%%%%%%%%%%%%%%
%%%%%%%%%%%%%%%%%%%%%%%%%%%%%%%%%%%%%%%%%%%%%%%%%%%%%%%%%%%
%%%%%%%%%%%%%%%%%%%%%%%%%%%%%%%%%%%%%%%%%%%%%%%%%%%%%%%%%%%
\TextSlide
{Allgemeine Schocksymptomatik}
{\label{topic:schock-allgemeine-symptomatik}%
\index{Schock!allgemeine Symptomatik}%
Die Symptome eines Schocks sind abhängig vom jeweiligen
Stadium. Problematisch ist, dass am Anfang die Symptome
eher diskret sind. Das erste Anzeichen eines beginnenden Schocks ist
oft lediglich eine kühle
Haut \cite{Kaplan2001}! Natürlich kann es
auch andere Ursachen für eine kühle Haut geben, z.\,B. eine
kalte Umgebungstemperatur. Die Tabelle
\FullPrettyRef{tab:schocksymptomatik}, listet typische
allgemeine Schocksymptome auf.


}
{%
\begin{itemize}
  \item Stadienabhängig
  \item \FullPrettyRef{tab:schocksymptomatik}
\end{itemize}
}
{}
{}
{}

\begin{Synopsis}
  \SynopsisItem{Ob ein Schock oder eine Schockgefahr besteht wird anhand von
  \underline{Schocksymptomen} und anhand der
  \underline{Verdachtsdiagnose} beurteilt.}
\end{Synopsis}

\begin{table}[H]
% \begin{WidePar}
\Captionof{table}{Allgemeine
Schocksymptomatik.\label{tab:schocksymptomatik}}
\CorrectionItemizeBox\FontTableBase

\begin{tabularx}{\linewidth}{p{9em}X}
\addlinespace\toprule\addlinespace
\noindent\TabHeading{Stadium}
&
\noindent\TabHeading{Symptome}
\\\addlinespace\midrule\addlinespace
\TabSub{\color{DarkGreen} \index{Schockgefahr}Schockgefahr}

\bigskip

\centering\includegraphics[width=1cm]{./icons/sm-basic.png}
&
 \begin{itemize}
     \item Keine eindeutigen Anzeichen
     \item Vorausschauend handeln!
     \item Verdachtsdiagnose bedenken.  Muss man damit
     rechnen, dass sich ein Schock entwickelt?  Besteht
     Schockgefahr?
 \end{itemize}
\\\addlinespace
\TabSub{\color{OrangeRed}Schock}

\bigskip

\centering\includegraphics[width=1cm]{./icons/sm-pale.png}
&
 \begin{itemize}
     \item HF ${\uparrow}$
     \item Mäßiger RR-Abfall, aber RR noch {\textgreater} HF
     \item Neurologisch unauffällig, ev. euphorisch oder
     agitiert
     \item Haut blass, evtl. kaltschweißig
     \item Venen sind kollabiert oder gestaut
     \item Peripherie schlecht durchblutet, Rekapzeit erhöht
 \end{itemize}
\\\addlinespace
\TabSub{\color{red}Schwerer Schock

Vollbild}

\bigskip

\centering\includegraphics[width=1cm]{./icons/sm-pale-frown.png}
&
\begin{itemize}
  \item Tachykardie (HF ${\uparrow}$, \enquote{fliehender} Puls)
  \item RR-Abfall (schwacher Puls, peripher schlecht oder nicht tastbar)
  \begin{itemize}
    \item Herzfrequenz {\textgreater} systolischer RR
  \end{itemize}
  \item Tachypnoe (schnelle, flache Atmung)
  \item Neurologisch auffällig, getrübt
  \item Haut blass, zyanotisch (bläulich), feucht
  \item Venen kollabiert oder gestaut
  \item Nur mehr schwer behebbar! ${\rightarrow}$ Schock  möglichst
  frühzeitig erkennen \&  behandeln
\end{itemize}
\\\addlinespace
\TabSub{\color{LightSlateGray}Endstadium}

\bigskip

% \huge\centering \bomb
&
\begin{itemize}
  \item Haut grau, klebrig
  \item Versagen der Herzfunktion
  \item RR kaum mehr messbar
  \item Evtl. Koma
  \item Evtl. Schnappatmung
  \item Patient ist sterbend.
\end{itemize}
\\\addlinespace\bottomrule
\end{tabularx}
% \end{WidePar}
\end{table}


% \clearpage % TEMP temp clearpage
\TextSlide{Schockschere}
{
\index{Schockschere}
\label{topic:schockschere}
Normalerweise ist die Herzfrequenz niedriger als der
systolische Blutdruck. (Achtung: \EE{Kinder haben andere
Normalwerte für HF und RR!} Die Faustregel ist deshalb nicht
anwendbar!) Wenn der Körper versucht den Schock zu
kompensieren, kann sich dieses Verhältnis umkehren. Man
spricht dann von der Schockschere, der Patient befindet sich
dann \EE{bereits  im Vollbild eines Schocks!} \E{Die
Schockschere ist ein \EE{spätes Zeichen!} Auch wenn sich ein
Patient (noch) nicht in der Schockschere befindet, kann er
trotzdem einen Schock haben!}



% % TODO Schockschere
% \Captionof{figure}{Schockschere. \AuthNotes{Achtung! Bild
% fehlt!}}


}{
\begin{itemize}
  \item RR\textsubscript{syst}\,{\textless}\,HF
  \item Dann bereits \E{schwerer} Schock
  \item Spätzeichen, auch ohne Schockschere kann Schock bestehen.
  \item[\Ca{!}] \EE{Schockbeurteilung}: Schocksymptome und
  Verdachtsdiagnose, \E{nicht} primär Schockschere!
  \item Für Kinder nicht anwendbar
\end{itemize}
}{}{}{}

\begin{Synopsis}
\SynopsisItem{Wenn sich das Verhältnis \enquote*{Herzfrequenz zu
systolischem Blutdruck} umkehrt, befindet sich der Patient
bereits in einem schweren Schock.} 
\SynopsisItem{Die Schockschere
ist ein \EE{Spätzeichen!} Auch wenn die Herzfrequenz
niedriger ist als der systolische Blutdruck kann sich der
Patient bereits im Schock befinden!} 
\SynopsisItem{Ob ein Schock
oder eine Schockgefahr besteht, wird anhand von
\E{Schocksymptomen} und anhand der
\E{Verdachtsdiagnose} beurteilt, \EE{\E{nicht}} primär anhand der Schockschere!}
\end{Synopsis}

\Text{Schock beim Kind}
{Die Schockschere gilt beim Kind nicht, da die
Normalwerte andere sind! Kinder bleiben außerdem lange
kreislaufstabil, verfallen dann aber schlagartig!

}{}
{}
{}
{}
{}

%%%%%%%%%%%%%%%%%%%%%%%%%%%%%%%%%%%%%%%%%%%%%%%%%%%%%%%%%%%
%%%%%%%%%%%%%%%%%%%%%%%%%%%%%%%%%%%%%%%%%%%%%%%%%%%%%%%%%%%
%%%%%%%%%%%%%%%%%%%%%%%%%%%%%%%%%%%%%%%%%%%%%%%%%%%%%%%%%%%

\Text
{Unterschiedliche Schockarten}
{\index{Schock!Arten}
\label{topic:schock-arten}%
Beim Schock können einzelne oder mehrere Teile des
Kreislaufsystems  (Blut, Gefäße, Herz) gestört sein.
Dementsprechend werden unterschiedliche Mechanismen
schlagend und man kann verschiedene Schockarten
unterscheiden, siehe
\FullPrettyRef{tab:schock-mechanismen}, und
\FullPrettyRef{tab:uebersicht-schockarten}.

}
{}
{}
{}
{}

\begin{table}[H]\CorrectionItemizeBox\FontTableBase
\caption{Mechanismen der Schockarten}
\label{tab:schock-mechanismen}
\begin{tabularx}{\linewidth}{XXX}
\toprule\addlinespace
\multicolumn{3}{c}{\noindent\TabHeading{Kreislauf}}
\\\addlinespace\midrule\addlinespace
\noindent\TabSub{Volumen} & \noindent\TabSub{Herz} & \noindent\TabSub{Gefäße} 
\\\addlinespace\cmidrule(lr){1-1}\cmidrule(lr){2-2}\cmidrule(lr){3-3}\addlinespace
Es ist \E{zu wenig Flüssigkeit} bzw. Blut
im Gefäßsystem vorhanden
&
Der \E{Herzauswurf} (bei normalem Volumen) ist
zu gering.
&
Die Gefäße sind inadäquat \E{weitgestellt},
das Blut versackt
\\\addlinespace
\begin{itemize}
  \item \EE{Hypovolämischer} Schock
\end{itemize}
&
\begin{itemize}
  \item \EE{Kardiogener} Schock
\end{itemize}
&
\begin{itemize}
  \item \EE{Anaphylaktischer} Schock
  \item \EE{Septischer} bzw. \EE{toxischer} Schock
  \item \EE{Neurogener} Schock
\end{itemize}
\\\addlinespace\bottomrule
\end{tabularx}
\end{table}



\Layout{57}
{}
{}
{}
{%
\small\FontSansNarrowFamily\linespread{1}\selectfont
\begin{tabularx}{\linewidth}{XXX}
\toprule\addlinespace
\noindent\TabHeading{Schockart} & \noindent\TabHeading{Ursachen} & \noindent\TabHeading{Spezielle Symptome}
\\\addlinespace\midrule\addlinespace
\TabSub{Hypovoläm}		

\DefinitionInPara{Verlust von Blutvolumen (Vollblut oder
Blutbestandteilen) aus den Gefäßen} 
&
\begin{itemize}
  \item Verlust nach außen
  \begin{compactitem}
    \item Blutung
    \item Verbrennungen (Hautbarriere versagt groß\-flächig)
    \item Gastrointestinal (Durchfall, Diarrhoe)
    \item Exsikkose (\FullPrettyRef{topic:exsikkose})
  \end{compactitem}
  \item Verlust nach innen
  \begin{compactitem}
    \item Innere Blutungen
    \begin{compactitem}
      \item Nach Trauma, rupturierte
      Eileiter\-schwanger\-schaft, \ldots
    \end{compactitem}
    \item Wassereinlagerung ins Gewebe
    \item \Z{Aszites
    (\enquote*{Wasserbauch}, z.\,B. bei
    Lebererkrankungen)}
  \end{compactitem}
\end{itemize}
& 
\begin{itemize}
  \item \FullPrettyRef{tab:schocksymptomatik}
\end{itemize}
\\\addlinespace\midrule\addlinespace
\TabSub{Kardiogen}	

\DefinitionInPara{Mangelnder Herzauswurf durch insuffiziente
Pumpleistung} 
&
\begin{itemize}
  \item Herzversagen, Herzinsuffizienz
  \begin{compactitem}
    \item Myokardinfarkt (häufigste Ursache)
    \item Herzrhythmusstörungen
  \end{compactitem}
  \item Lungenembolie
\end{itemize}
&
\begin{itemize}
    \item RR-Abfall
    \item Stauungszeichen (Lungenödem, Beinödeme, gestaute Halsvenen)
    \item Atemnot
  \end{itemize}
\\\addlinespace\midrule\addlinespace
\TabSub{Anaphylaktisch}

\DefinitionInPara{Überschießende
allergische Reaktion}
&
\begin{itemize}
  \item Überschießende
  Reaktion der Körperabwehr
  \item Gefäße erweitert, Blut versackt
  \item Durchlässigkeit der Gefäße ↑, Plasmaverlust in
  Zwischengewebsraum
  \item Mögliche Auslöser:
  \begin{compactitem}
    \item Pollen
    \item Medikamente
    \item Kontrastmittel
    \item tierische Gifte
    \item u.\,v.\,a.\,m.
  \end{compactitem}
\end{itemize}
&
\begin{itemize}
%   \item \FullPrettyRef{tab:anaphylaxie-stadien}
  \item Stadienabhängig
  \item[\TabSub{I}]
  Hautreaktion: evtl. Ausschlag, Rötung;
  {Quaddeln}, Flush, Ödeme
  
  Allgemeinsymptome: Juckreiz, Zittern, Schwindel, Kopfschmerz
  \item[\TabSub{II}]
  RR ${\downarrow}$, Tachykardie
  
  Übelkeit, Erbrechen, Durchfall
  \item[\TabSub{III}]
  Allgemeine Schocksymptome
  
  Bronchospasmus
  \item[\TabSub{IV}]
  Atem-/Kreislaufstillstand
  
  Schwerer Schock
\end{itemize}
\\\addlinespace\midrule\addlinespace
\TabSub{Septisch, toxisch}

\DefinitionInPara{Schwere
körperweite Infektion oder Einwirkung eines die Gefäße
beeinflussenden Giftstoffs (Toxin)}
&
\begin{itemize}
  \item Langanhaltender Entzündungsherd der plötzlich Erreger
  \enquote*{ins Blut streut}
  \item Fremdkörper (Venenverweilkanülen, Tampons, \ldots)
  \item \ldots
\end{itemize}
&
\begin{itemize}
  \item Schock
  \item Hohes oder kein Fieber
  \item Haut heiß, trocken, rot, später bleich, grau,
  marmoriert
\end{itemize}
\\\addlinespace\midrule\addlinespace
\TabSub{Neurogen}

\DefinitionInPara{Unterbrechung oder Störung der nervlichen
Versorgung der Kreislaufregulation. Es kommt zu inadäquater arterieller
und venöser Gefäßerweiterung \Z{(Vasodilatation)}.}

&
\noindent\begin{itemize}
    \item Rückenmarksverletzung (Querschnitt-Syndrom)
    \item  Gefäßwandmuskulatur bekommen keine Nerven-Impulse
    und kontrahiert sich nicht mehr
    \item Gefäße weiten sich ${\rightarrow}$ Blut versackt
\end{itemize}
&
\begin{itemize}
  \item \FullPrettyRef{tab:schocksymptomatik}
\end{itemize}
\\\addlinespace\bottomrule\addlinespace
\end{tabularx}
}
{Übersicht: \SlideHeading{Schockarten}
\label{tab:uebersicht-schockarten}}
{}




\Text
{Hypovolämischer Schock}
{\label{topic:schock-hypovolaemischer}
\index{Schock!Hypovolämischer}
\DefinitionInPara{Schock aufgrund eines Verlustes von
Flüssigkeit aus den Gefäßen.}
Es können dabei 
(Voll-) Blut, oder auch nur Blutbestandteile wie Plasma oder
Wasser und Elektrolyte
verloren gehen. Der Verlust kann jeweils nach
innen und/oder nach außen erfolgen, siehe
\FullPrettyRef{tab:uebersicht-urachen-fluessigkeitsverlust-schock}.
Bei akutem Flüssigkeitsverlust kann ein Verlust von bis zu
20\PC* des Blutvolumens  gut kompensiert werden. Auf größere
Verluste  folgt ein Absinken von Herzauswurf und RR!}
{}
{}
{}
{}




\begin{table}
\caption{Übersicht: Ursachen für Flüssigkeitsverluste beim
Schock.}
\label{tab:uebersicht-urachen-fluessigkeitsverlust-schock}
\CorrectionItemizeBox\FontTableBase
\begin{tabularx}{\linewidth}[H]{XX}
\addlinespace\toprule\addlinespace
Flüssigkeitsverlust nach außen
&
Flüssigkeitsverlust nach innen
\\\addlinespace\midrule\addlinespace
\begin{itemize}
  \item Blutung
  \item Verbrennungen (Hautbarriere versagt  groß\-flächig)
  \item Gastrointestinal (Durchfall, Diarrhoe)
  \item Harnausscheidung (Diabetes mellitus,
  Diabetes insipidus, Diuretika, Polyurie bei
  Nierenversagen)
\end{itemize}
&
\begin{itemize}
  \item Innere Blutungen
  \begin{itemize}
    \item Nach Trauma, rupturierte Eileiter\-schwanger\-schaft, etc
  \end{itemize}
  \item Wassereinlagerung ins Gewebe
  \item \Z{Aszites
  (\enquote*{Wasserbauch}, z.\,B. bei
  Lebererkrankungen)}
\end{itemize}
\\\addlinespace\bottomrule
\end{tabularx}
\end{table}



\AuthNotes{Begriff \enquote{Schocklagerung}
ein für alle mal in der Ausbildung streichen, nur Hinweis
dass dieser Ausdruck veraltet ist sollten sie ihn wo
aufschnappen. Stattdessen \enquote{Beine-hoch}-Lagerung.

Wenn KI gegen Beine hoch: Flachlagerung bzw. stab. Seitenlage empfehlen?

STEU: ja, Flachlagerung wenn KI. 2009-06-28}

\TextSlide
{Kardiogener Schock}
{\index{Schock!Kardiogener}
\label{topic:schock-kardiogener}
\DefinitionInPara{Schock infolge mangelndem Herzauswurfes
durch insuffiziente Pumpleistung.} 
Zu einem kardiogenen Schock kann es bei einem \EE{akut
auftretenden Herzversagen}, z.\,B. durch einem akuten
\E{Myokardinfarkt} (häufigste Ursache) oder
\E{Herzrhythmusstörungen}, kommen. Die \E{Lungenembolie} als
Auslöser eines Schocks wird auch zum kardiogenen Schock
gerechnet, da es durch die Verstopfung des kleinen
Kreislaufes zu einer mangelhaften Füllung des Herzens und
damit zu einer mangelhaften Auswurfleistung kommt.
Typisch ist ein Blutdruckabfall. Durch den
\EE{Blutrückstau} kann es zu Stauungszeichen kommen
(Lungenödem bei Linksherzversagen, Beinödeme und gestaute
Halsvenen bei Rechtsherzversagen). Häufig klagen die
Patienten über \EE{Atemnot}.
}
{%
\FullPrettyRef{tab:uebersicht-schockarten}
}
{}
{}
{}



\TextSlide
{Anaphylaktischer Schock}
{\label{topic:schock-anaphylaktischer}
\index{Schock!Anaphylaktischer}
\index{Schock!allergischer}
\DefinitionInPara{Schock aufgrund einer überschießenden
allergischen Reaktion -- \enquote{Allergischer
Schock}.}
Ein anaphylaktischer Schock wird durch eine
über\-schieß\-ende Reaktion der Körperabwehr (Immunreaktion)
auf bestimmte Stoffe (Antigene) ausgelöst.
Durch von den Immunzellen gebildete Stoffe werden
die Gefäße erweitert, zusätzlich steigt die
Durchlässigkeit der Gefäße\ZFootnote{(Gefäßpermeabilität)}.
Das Blut versackt und es kommt zu einem  Plasmaverlust in
den Zwischengewebsraum\ZFootnote{(Interstitium)}; der
Blutdruck fällt massiv ab!
Mögliche \E{Auslöser} können sein: 
\begin{compactitem}
    \item Pollen
    \item Medikamente, Antibiotika
    \item Kontrastmittel (für radiologische Untersuchungen)
    \item tierische Gifte (Bienenstich etc.)
    \item Lebensmittel (Meeresfrüchte, Nüsse, etc.)
    \item und vieles andere \ldots
\end{compactitem}

Die {\TxSymptome} sind Abhängig vom Schweregrad und Stadium,
siehe dazu \FullPrettyRef{tab:anaphylaxie-stadien}.
}
{%
\begin{itemize}
  \item \FullPrettyRef{tab:uebersicht-schockarten}
\item {\TxSymptome}: \FullPrettyRef{tab:anaphylaxie-stadien}
\end{itemize}
}
{}
{}
{}




\begin{table}[H]
\CorrectionItemizeBox\FontTableBase
\Captionof{table}{Stadien der
Anaphylaxie\label{tab:anaphylaxie-stadien}}
\begin{tabulary}{\linewidth}[H]{cL}\addlinespace\toprule\addlinespace
\noindent\TabHeading{\mbox{Stadium}}
 &
\noindent\TabHeading{{Symptome}}
\\\addlinespace\midrule\addlinespace
\TabSub{I}
 &
Hautreaktion: evtl. Ausschlag, Rötung;
{Quaddeln}, Flush, Ödeme

Allgemeinsymptome: Juckreiz, Zittern, Schwindel, Kopfschmerz

\\\addlinespace
\TabSub{II}

 &
RR ${\downarrow}$, Tachykardie

Übelkeit, Erbrechen, Durchfall 

\\\addlinespace
\TabSub{III}
 &
Allgemeine Schocksymptome

Bronchospasmus
\\\addlinespace
\TabSub{IV}
 &
Atem-/Kreislaufstillstand

Schwerer Schock
\\\addlinespace\bottomrule
\end{tabulary}
\end{table}

% \MassnahmenNG{Spezielle Maßnahmen}{Anaphylaktischer
% Schock}{}{%
% \begin{itemize}
%   \item Allgemein: \EE{Allgemeine Maßnahmen der
%   Schockbekämpfung}
%   (\prettyref{m:am-schock})
%   \item Antigen  nach Möglichkeit entfernen!
%   \item Lagerung: situationsgerecht.
%   Grundsätzlich Beine hoch, bei
%   Atemproblemen jedoch stark erhöhten
%   Oberkörper erwägen.
% \end{itemize}
% }

\TextSlide
{Septischer Schock und Toxischer Schock}
{%
\index{Schock!septischer}
\index{Schock!toxischer}
\label{topic:schock-septischer}
\label{topic:schock-toxischer}
\DefinitionInPara{Schock infolge einer schweren,
körperweiten Infektion oder aufgrund eines, 
die Gefäße beeinflussenden, Giftstoffes
(Toxins).}
Siehe auch \FullPrettyRef{topic:sepsis}.
\index{Toxine!Toxischer Schock}%
Ein septisch-/toxischer
Schock wird durch diverse Bakterien und Toxine verursacht.
Die Erreger können auf unterschiedliche Wege in die Blutbahn
gelangen. Oft liegt ein chronischer Entzündungsherd vor, von
dem aus Keime in die Blutbahn streuen.  Mit Keimen
besiedelte und im Körper gelegene Fremdkörper, darunter
fallen auch medizinische Materialien wie
lang liegende Venenverweilkanülen, Harnkatheter u.\,ä., aber
auch Hygieneartikel wie vergessene Tampons, sind ideale
Brutstätten, da hier das Immunsystem des Körpers keinen Zugang hat. 

Durch die bakterielle oder toxische Gefäßwandschädigung kommt es
zum Flüssigkeitsaustritt in den Zwischengewebsraum sowie zu Gefäßweitstellung.
Man kann die allgemeinen Schocksymptome wahrnehmen. Der
Patient kann sowohl hohes  oder auch kein Fieber haben
Die Haut ist anfänglich heiß, trocken und rötlich, später
ist sie eher bleich, grau oder sogar marmoriert.
}
{\FullPrettyRef{tab:uebersicht-schockarten}}
{}
{}
{}



\TextSlide
{Neurogener Schock}
{%
\index{Schock!Neurogener}
\index{Schock!Spinaler}
\DefinitionInPara{Schockzustand ausgelöst durch eine
Unterbrechung oder Störung der nervlichen Versorgung der
Kreislaufregulation. Es kommt zu inadäquater arterieller
und venöser Gefäßerweiterung.
\Z{(Vasodilatation)}}
Auch die Muskulatur der Blutgefäßwand wird mit Nerven
versogrt, damit die Weite der Gefäße reguliert werden kann.
Im Rahmen von Erkrankungen und Verletzungen kann diese
nervale Versorgung unterbochen sein, in weiterer Folge
weiten sich die betroffenen Gefäße und das Blut versackt. 
Typischerweise kommt dies im Rahmen einer
\EE{Wirbelsäulenverletzung} vor, bei der das Rückenmark
geschädigt bzw. unterbrochen wird (\I{Querschnitt-Syndrom}).

}
{\FullPrettyRef{tab:uebersicht-schockarten}}
{}
{}
{}

\label{topic:schockbekaempfung}
\begin{Synopsis}
  \SynopsisItem{Grundsatz: \EE{Der Schock muss bekämpft werden
  bevor man ihn bemerkt.}}
\end{Synopsis}



% M %%%%%%%%%%%%%%%%%%%%%%%%%%%%%%%%%%%%%%%%%%%%%%%%%%% M %
\ProcedurePrintDefault{MR57XX0B}


%%%%%%%%%%%%%%%%%%%%%%%%%%%%%%%%%%%%%%%%%%%%%%%%%%%%%%%%%%%
%%%%%%%%%%%%%%%%%%%%%%%%%%%%%%%%%%%%%%%%%%%%%%%%%%%%%%%%%%%
%%%%%%%%%%%%%%%%%%%%%%%%%%%%%%%%%%%%%%%%%%%%%%%%%%%%%%%%%%%
\subsection{\HeadingKeyword{Kreislaufstillstand}}
\label{bkm:RefHeading42514602}
\label{topic:kreislaufstillstand}
\index{Kreislaufstillstand}


\Text
{\TxBeschreibung}
{%
\DefinitionInPara{Ein Kreislaufstillstand ist ein totales
Pumpversagen des Herzens und Stillstand des Kreislaufes.}
Das Pumpversagen kann unterschiedliche Ursachen haben:
\begin{itemize}
  \item Herzrhytmusstörungen
  \begin{itemize}
    \item Asystolie (keine Herztätigkeit)
    \item Kammerflimmern (unkoordinierte Herztätigkeit)
    \item Zu schnelle Herzfrequenz, bei der sich das Herz
    nicht mehr füllen oder etwas auswerfen kann ((Pulslose)
    Ventrikuläre Tachykardie)
    \item Elektromechanische
    Entkopplung (pulslose elektrische Aktivität): keine
    Pumpreaktion auf Signale des Reizleitungssystems im Herz
  \end{itemize}
  \item Ischämie
  \item Mechanische Blockaden
  \item Blutverlust
\end{itemize}
}
{}
{}
{}
{}

\Text{Ursachen}
{
\begin{itemize}
  \item Minderversorgung des Herzmuskels mit \ce{O2}\par
  \begin{inparaitem}
    \item Herzinfarkt,
    \item Herzinsuffizienz, \ldots
  \end{inparaitem}
  \item Störungen der Atmung\par
  \begin{inparaitem}
    \item Asthma bronchiale,
    \item Lungenembolie,
    \item zentrale Atemregulationsstörungen, \ldots
  \end{inparaitem}
  \item Sonstige Störungen\par
  \begin{inparaitem}
    \item Trauma,
    \item Vergiftungen,
    \item thermische Schäden,
    \item Elektrolytverschiebungen, \ldots
  \end{inparaitem}
\end{itemize}
}{}{}{}{}

\TextSlide{Folgen und Symptome}
{Unmittelbar nach Aussetzen des Herzschlages ist der Patient
\E{pulslos} und der Blutkreislauf steht still. Es wird ihm
schwindlig und schlecht, nach kurzer Zeit (ca.
\VonBis{15}{20}\Sek*) wird er \E{bewusstlos}. Nach weiteren
10\Sek* tritt ein \E{Atemstillstand} ein (oder eine
\I{Schnappatmung}).
Nach \VonBis{3}{5}\Min* treten Hirnschäden auf, welche nicht
mehr reversibel sind. In kühler Umgebung kann diese
Zeitspanne etwas ausgedehnt werden.
Viele Patienten haben, auch nach erfolgreicher Reanimation,
im weiteren Verlauf neurologische Defizite.
}{
\begin{itemize}
  \item Aussetzen des Herzschlages $\rightarrow$ 
  sofort pulslos!
  \item[\Sign{→}] Bewusstlosigkeit nach \VonBis{15}{20}\Sek*
  \item[\Sign{→}] Atemstillstand/Schnappatmung nach ca. 10\Sek*
  \item[\Sign{→}] \EE{Irreversible Hirnschäden nach \VonBis{3}{5}\Min*}
\end{itemize}

}{}{}{}




\TextSlide
{\TxQuerverweise}
{\begin{itemize}
  \item Reanimation, Erste Hilfe:
  \FullPrettyRef{topic:reanimation-eh}
  \item Basic Life Support:
  \CO{[BLS]}, \FullPrettyRef{chp:BLS}
  \item Advanced Life Support: \CO{[ALS]},
  \FullPrettyRef{chp:ALS}
\end{itemize}}
{}
{}
{}
{}



\subsection{Störungen des \HeadingKeyword{Blutdrucks}}
\subsubsection{\HeadingKeyword{Hypotonie}}
\label{topic:hypotonie}


 
\TextSlide{{\TxBeschreibung} und Richtwert}
{%
\DefinitionInPara{Als \T{Hypotonie} bezeichnet man einen 
unangemessen zu niedrigen Blutdruck.}
Beim Erwachsenen beträgt der \E{untere} Grenzwert des
systolischen Blutdrucks 100\MmHg*. 
Dieser Grenzwert ist jedoch sehr individuell. Viele
Menschen haben auch unter Normalbedingungen einen
systolischen Blutdruck von 90\MmHg* oder sogar noch
weniger, ohne Beschwerden zu haben. Dies muss bei der
Interpretation des Wertes berücksichtigt werden. Bei Kindern
gelten grundsätzlich andere Nomalwerte.
\citep{PraxisleitAllgemeinmed6} }{%
\begin{itemize}
  \item Erwachsener: RR\textsubscript{sys} {\textless}
  ca. 100\MmHg*
  \item Individuell unterschiedlich
  \item Kinder haben andere Normalwerte
\end{itemize}
}{}{}{}


\subsubsection{\HeadingKeyword{Kollaps} und \HeadingKeyword{Synkope}}
\label{topic:kollaps}
\label{topic:synkope}
\index{Kollaps|TLike}
\index{Synkope|TLike}


\Text{{\TxBeschreibung} und Ursachen}
{%
\DefinitionInPara{Als Kollaps bzw. Synkope bezeichnet man 
eine kurzfristige Kreislaufregulationsstörung.} 
Manche Autoren unterscheiden zwischen beiden Begriffen und 
ordnen der Synkope eher kardiale Ursachen zu. 

Ein Kollaps kann seine Ursachen in verschiedenen
Organssystemen haben:
\begin{itemize}
  \item Herz (z.\,B. Herzrhythmusstörungen \Z{(\I{Adams-Stokes-Anfall})})
  \item Gefäßen (z.\,B. Hitzebedingte Erweiterung
  der Gefäße mit \enquote*{Versacken} des Blutes in den
  Beinen)
  \item Blut (z.\,B. Blutarmut (\I{Anämie}))
  \item Hirn
  \item Hirngefäßen (z.\,B. TIA,
  \prettyref{topic:tia} ) \ldots
  \item Sonstige Ursachen und Faktoren: Hypoglykämie, 
  Menstruation
\end{itemize}

Eine Abklärung ist präklinisch oft nicht möglich, und auch
unter stationären Bedingungen oft langwierig und nicht
selten ergebnislos. Da jedoch schwerwiegende und
lebensbedrohliche Diagnosen möglich sind, muss bei unklaren
Kollapszuständen eine gründliche Abklärung erfolgen.

Es gibt jedoch auch typische \enquote*{Kollaps-Situationen}:
Ein schwüler Tag in einer überfüllten U-Bahn-Garnitur, nicht
gefrühstückt, auf dem Weg zu einem Vorstellungstermin und
seit gestern in der Menstruation. Oder im stickigen
Veranstaltungszelt bei 40\Deg , wobei der der ganze
Flüssigkeitskonsum aus 5 (entwässernden) Dosen
RedBull{\texttrademark}
bestand. So mancher war auch ob eines Begräbnisses so
ergriffen, dass er sich fast in die Grube neben den Sarg
dazugelegt hätte \ldots 
}{%

}{}{}{}


\TextSlide{\TxSymptome}
{%
Der Patient berichtet über Schwindel, evtl. ist er
zusammengesackt, war kurz bewusstlos. Der Patient wacht oft
auf dem Boden liegend auf und weiß nicht, was geschehen ist.
Er erzählt, es wäre ihm schwarz vor den Augen geworden. Der
Blutdruck ist meist niedrig, der Patient ist eventuell
bradykard oder tachykard.\bigskip
}{
\begin{compactitem}
  \item Schwindel
  \item Evtl. kurze Ohnmacht, Schwarzwerden vor Augen
  \item RR ${\downarrow}$, evtl. Bradykardie oder Tachykardie
\end{compactitem}
}{}{}{}





% M %%%%%%%%%%%%%%%%%%%%%%%%%%%%%%%%%%%%%%%%%%%%%%%%%%% M %
% Kollaps/Synkope
\ProcedurePrintDefault{MR55XX0C}


\subsubsection{Arterielle \HeadingKeyword{Hypertonie}}
\label{topic:hypertonie}
\label{topic:arterielle-hypertonie}
\index{Hypertonie!arterielle}



\TextSlide{{\TxBeschreibung}}
{%
\DefinitionInPara{Die arterielle \T{Hypertonie}
(\I{Bluthochdruck}, \I{Hochdruckkrankheit}) ist eine meist
symptomlose, \textit{chronische} Erkrankung, bei welcher der
arterielle Blutdruck dauerhaft erhöht ist.} Sie kann auf
Dauer zu schweren Schäden an den Blutgefäße und im Zuge 
von plötzlichen \E{Blutdruckkrisen} auch zu akuten
Problemen führen. Aufgrund der Schädigung der Blutgefäße
ist sie ein ein wesentlicher \E{Risikofaktor} für
Erkrankungen wie die koronare Herzkrankheit, Schlaganfälle
etc.\Commentary[Hypertonie]{Neben der systemischen
arteriellen (essentiellen) Hypertonie gibt es auch andere 
Krankheitsbilder, bei denen der Druck in anderen
Blutgefäßen erhöht ist, z.\,B. die pulmonale Hypertonie.} 

Bei der chronischen Hypertonie liegt der obere
Grenzwert im Allgemeinen\ZFootnote{Grenzwert Blutdruck: 
Bei Patienten mit Risikofaktoren gelten andere Grenzwerte.} 
beim Erwachsenen bei \EE{\RRmmHg{140}{90}}. Dieser Grenzwert
ist nur für die dauerhafte Behandlung maßgeblich, akute
Beschwerden sind erst bei deutlich höheren Werten zu
erwarten. Eine Ausnahme stellen \EE{Schwangere} dar:
Eine Hypertonie kann ein Zeichen einer schwerwiegenden
Schwangerschaftserkrankung sein und muss zeitnah
entsprechend abgeklärt und behandelt werden.
}
{%
\begin{itemize}
  \item Chronische Erhöhung des arteriellen Blutdrucks
  \item Risikofaktor für Herz-/Kreis\-lauf\-er\-krank\-ung\-en
\end{itemize}
}


% \TextSlide{Ursachen}
% {
% \begin{itemize}
%   \item Widerstandserhöhung (durch Engstellung der Gefäße)
%   \item Vermehrte Herzarbeit
%   \item Sympathikusaktivierung, Adrenalin und Cortisol
%   (Stresshormone) durch Dauerstress,  Bewegungsmangel, genetische Veranlagung
%   \item u.\,a. \ldots
% \end{itemize}
% 
% }
% {
% 
% }{}{}{}

\begin{Danger}
  \item Bei \EE{Schwangeren} sind die Grenzwerte genau zu beachten!
  (\EE{EPH-Gestose}, \prettyref{topic:eph-gestose})!
\end{Danger}

\Layout{60}{Langzeitfolgen}{%
Die Hypertonie hat in der Notfallmedizin auch eine wichtige
indirekte Bedeutung: Sie ist Ursache für viele andere
Erkrankungen, die ihrerseits akut lebensbedrohlich werden
können.
Zwei Wirkungen der Hypertonie sind dabei besonders wichtig:
Die \EE{Gefäßschädigung} und die \EE{Belastung des Herzens
durch den hohen Gegendruck}.

Bei der Gefäßschädigung kommt es zur Ablagerung durch Kalk
und sonstige Plaques und über längere Zeit zur \E{Verengung
oder Verstopfung der Gefäße}. Je nach Organ kann das sehr
massive Auswirkungen haben. Im Herz führt es zur \E{Angina
Pectoris} oder zum \E{Herzinfarkt}, im Hirn zu
\E{Schlaganfällen}, in den Nieren zu einer
\E{Niereninsuffizienz}, usw. Selbst große Gefäße wie die
Hauptschlagader (Aorta) können derartig geschädigt werden,
dass sie reißen können (Aortenaneurysma) -- eine für den
Patienten sehr lebensgefährliche Komplikation.

Die Schädigung des Herzmuskels ist die Folge von dem
erhöhten Druck gegen den das Herz pumpen muss. Der Muskel
vergrößert sich um mit der Belastung fertig zu werden, kann
aber dann nicht mehr ausreichend mit sauerstoffreichen Blut
versorgt werden.

Die Hypertonie ist somit ein wichtiger
\E{Risikofaktor} für viele schwere Notfälle. Daher ist
es auch verständlich, dass die Behandlung der Hypertonie
einen großen Stellenwert in der medizinischen Versorgung der
Bevölkerung hat.
}{
\begin{itemize}
  \item \EE{Gefäßschädigung}
  \begin{itemize}
    \item Herzkranzgefäße ${\rightarrow}$ Infarkt
    \item Hirngefäße ${\rightarrow}$ Schlaganfall
    \item Nierengefäße ${\rightarrow}$ Niereninsuffizienz ${\rightarrow}$
    Dialyse ${\rightarrow}$ typischer KTW-Patient
    \item Hauptschlagader ${\rightarrow}$ Aneurysma
  \end{itemize}
  \item \EE{Herzmuskelschädigung} \par  (Herz muss gegen
  größeren Widerstand pumpen)
  \begin{itemize}
    \item Krankhafte Vergrößerung des Herzmuskels
  \end{itemize}
\end{itemize}
}
{./images/teaser/imgp0481.jpg}
{}
{GABS}




\subsubsection{ \HeadingKeyword{Hypertensive Krise} und Hypertensiver \HeadingKeyword{Notfall}}
\label{topic:hypertensive-krise}
\label{topic:hypertensiver-notfall}

\Text{\TxBeschreibung}
{%
\DefinitionInPara{Eine plötzliche, sehr starke Erhöhung des
Blutdrucks wird als \T[Krise!hypertensive]{hypertensive Krise} oder
\T[Notfall!hypertensiver]{hypertensiver Notfall} bezeichnet.
Der Unterschied zwischen den beiden Diagnosen ist das Fehlen bzw. 
Vorhandensein von typischen Symptomen 
einer Endorganstörung (z.\,B. Kopfschmerzen, Sehstörungen, 
Brustschmerzen, Übelkeit,
Erbrechen, Schwindel \ldots).}
Vgl.
\prettyref{tab:hxpertension-krise-notfall}.

}{}{}{}{}

% \TextSlide{\TxSymptome}
% {%
% Vgl. \FullPrettyRef{tab:hxpertension-krise-notfall}.
% Ein erhöhter Blutdruck alleine ist oft symptomlos und wird
% vom betreffenden Patienten oft gar nicht bemerkt.
% Klassisch kann er sich durch \I{Nasenbluten}
% (\T{Epistaxis}\index{Epistaxis}) bemerkbar machen, welches
% ein typisches Symptom für einen plötzlich erhöhten Blutdruck
% ist.
% 
% Wenn Symptome wie Kopfschmerzen, Schwindel und Übelkeit,
% Brustschmerzen, Sehstörungen, etc. auftreten, ist dies ein
% Zeichen für die blutdruckbedingte Störung von Organen
% (\E{Organstörungen}). Diese Symptome einer Organstörung sind
% typische Leitsymptome für einen \E{hypertensiven
% \EE{Notfall}}.
% 
% }{
% \FullPrettyRef{tab:hxpertension-krise-notfall}.
% }{}{}{}

\Layout{57}{}{}{}{
\begin{tabularx}{\textwidth}{XX}
\toprule\addlinespace
\noindent\TabHeading{Hypertensive Krise}
&
\noindent\TabHeading{Hypertensiver Notfall}
\\\addlinespace\midrule\addlinespace
\Abk{RR}~{\textgreater}~\RRmmHg{230}{130} (Richtwert) 
bzw. deutliche individuelle Erhöhung

\EE{Patient fühlt sich gut}, evtl. Nasenbluten
 &
\Abk{RR}~{\textgreater}~\RRmmHg{230}{130} (Richtwert)
bzw. deutliche individuelle Erhöhung

\EE{+ Symptome einer Organstörung}
\\\addlinespace
Zufallsbefund? --
Oder steht der Blutdruck doch in Zusammenhang mit der
Berufung?

% \textit{Patient hat überhaupt keine Beschwerden?} Warum hat
% er Sie dann überhaupt gerufen?
% 
% Wenn es wirklich nur ein Zufallsbefund ist
% (\enquote*{Gesundheitstage},
% \enquote*{Vorsorgeuntersuchung}):
&
\begin{itemize}
    \item Kopfschmerz
    \item Sehstörungen
    \item Brustschmerz
    \item Übelkeit, Erbrechen, Schwindel
\end{itemize}
\\\addlinespace
Eine notfallmäßige Blutdrucksenkung ist (eher) nicht erforderlich.
&
\Ca{Vitale Bedrohung}: \EE{Gefahr des Herzversagens: Das
Herz muss plötzlich gegen großen Druck pumpen}

Eine frühzeitige Blutdrucksenkung 
ist notwendig. 
\\\addlinespace\bottomrule
\end{tabularx} 
}
{Unterscheidung Hypertensive Krise und Hypertensiver Notfall
\protect\citep{WHOISH2003,JNC7Report,BLIM3,IntroClinEmergMed4,Vlcek2008}.\label{tab:hxpertension-krise-notfall}}
{}

\begin{Synopsis}
  \SynopsisItem{Der Unterschied zwischen einer hypertensiven Krise
  und dem hypertensiven Notfall ist das Vorhandensein von
  Symptomen.}
\end{Synopsis}

\Commentary[Hypertensive Krise/Notfall]{(Die Unterscheidung
zwischen \enquote{Hypertensiver Krise} und \enquote{Hypertensiven Notfall} ist in der
Literatur nicht einheitlich. Oft wird auch nur
zwischen \enquote{Hypertensiver Krise ohne Symptome} bzw. \enquote{mit
Symptomen} unterschieden oder generell eine andere
Einteilung gewählt. Von diesen Spitzfindigkeiten abgesehen
gilt das Vorhandensein von Symptomen durchgehend als
entscheidendes Kriterium.
\citep{ClassenInnere5,KlinLeitAlle3,Fauci2008})}
% Renz-Polster2006 p442
% IntroClinEmergMed4 Kap. 28, p 393

\TextSlide
{\TxABCDE}
{~
\begin{ABCDEText}
  \item[\ABCDE{2}] Oft hat der Patient einen
  hochroten Kopf.
  \item[\ABCDE{4}] Ein erhöhter Blutdruck alleine ist oft symptomlos und wird
vom betreffenden Patienten oft gar nicht bemerkt. Häufig ist
\I{Nasenbluten} (\T{Epistaxis}\index{Epistaxis}) das
  erste Zeichen einer akuten Hypertonie. Beim
  hypertensiven \EE{Notfall} kommt es zu Symptomen für die
  blutdruckbedingte Störung von Organen
  (\E{Organstörungen}), z.\,B. oft zu {\RedFlag\,}neurologischen
  Symptomen wie {\RedFlag\,}Kopfschmerzen und {\RedFlag\,}Sehstörungen, weiters
  evtl. zu {\RedFlag\,}Brustschmerzen, Übelkeit,
  {\RedFlag\,}Erbrechen und {\RedFlag\,}Schwindel.
  \item[\ABCDE{B}] Evtl. kommt es aufgrund der kardialen Belastung zur {\RedFlag\,}Atemnot.
  \item[\ABCDE{C}] In sehr schweren Fällen, v.\,a. wenn eine
  kardiale Grunderkrankung besteht, können sich Symptome
  eines {\RedFlag\,}kardiogenen Schocks zeigen.
  \item[\ABCDE{D}] Es ist eine deutliche Blutdruckerhöhung
  (Richtwert \Abk{RR}~{\textgreater}~\RRmmHg{230}{130}
  bzw. deutliche individuelle Erhöhung) festzustellen
\end{ABCDEText}
}
{
\begin{ABCDESlide}
  \item[\ABCDE{2}] Hochroter Kopf 
  \item[\ABCDE{4}] Nasenbluten, {\RedFlag\,}Kopfschmerzen, {\RedFlag\,}Sehstörungen, 
  {\RedFlag\,}Brustschmerzen, Übelkeit, {\RedFlag\,}Erbrechen, {\RedFlag\,}Schwindel
  \item[\ABCDE{B}] Evtl. {\RedFlag\,}Dyspnoe
  \item[\ABCDE{C}] Evtl. {\RedFlag\,}kardiogener Shock
  \item[\ABCDE{D}] RR $\uparrow$ (Richtwert\Abk{RR}~{\textgreater}~\RRmmHg{230}{130} 
  bzw. deutliche individuelle Erhöhung)
\end{ABCDESlide}
}
{}
{}
{}

\TextSlide
{{\TxSAMPLER}}
{
\begin{SAMPLERText}
  \item[\SAMPLER{S}] Siehe ABCDE.
  \item[\SAMPLER{M}] Meistens ist bereits eine medikamentöse
  Einstellung durch den Hausarzt mittels eines oder
  Kombination mehrerer Medikamente
  erfolgt\ZFootnote{Gängige Antihypertensive Medikamente:
  \begin{inparaitem}
    \item Beta-Blocker: 
    Bisoprolol (Concor{\texttrademark}),
    Carvedilol (Dilatrend{\texttrademark}), 
    Metoprolol (Beloc{\texttrademark}, Seloken{\texttrademark}),
    Nebivolol (Nebilan{\texttrademark}, Nomexor{\texttrademark})
    \item ACE-Hemmer: 
    Enalapril (Mepril{\texttrademark}, Renitec{\texttrademark}), 
    Lisinopril (Acemin{\texttrademark}, Acetan{\texttrademark}) 
    Ramipril (Tritace{\texttrademark}),
    Fosinopril (Fositens{\texttrademark}), 
    Captopril
    \item AT-II-Inhibitoren: 
    Candesartan (Atacand{\texttrademark}, Blopress{\texttrademark}), 
    Losartan (Cosaar{\texttrademark}),
    Telmisartan (Micaardis{\texttrademark}),
    Valsartan (Diovan{\texttrademark})
    \item Kalzium-Antagonisten: 
    Diltiazem,
    Verapamil (Isoptin{\texttrademark}),
    Amlodipin (Norvasc{\texttrademark}),
    Lercanidipin (Zanidip{\texttrademark}),
    Nifedipin (Adalat{\texttrademark}, Buconif{\texttrademark})
    \item Entwässerung (Diuretika): 
    Furosemid (Lasix{\texttrademark}, Furon{\texttrademark}),
    Spironolacton (Aldactone{\texttrademark}, Spirobene{\texttrademark}), 
    Hydrochlorthiazid (HCT) üblicherweise in Kombination mit andren Medikamenten
    \item Kombinationspräparate: Diverse Kombinationen, oft
    mit Zusätzen wie \enquote*{Co-}, \enquote*{-comb},
    \enquote*{-comp}, \enquote*{plus} oder Zusammensetzung
    aus Wirkstoff- bzw. Markennamen. Z.\,B.:
    Acecomb{\texttrademark}, Exforge{\texttrademark},
    Fosicomb{\texttrademark},
  \end{inparaitem} }.
  Entscheidend ist die Frage, ob diese Medikation auch tatsächlich eingenommen wurde, 
  Fehler bei
  der Einnahme (vergessen, Packung aufgebraucht, \ldots)
  sind eine gute Erklärung für eine hypertensive Krise. 
  Ebenso kann es nach einer
  Medikamentenumstellung zur Blutdruckentgleisungen kommen. 
  \item[\SAMPLER{P}] Eine chronische arterielle Hypertonie
  ist oft schon bekannt und der Patient befindet sich
  deswegen bereits in Behandlung. Oft führt der Patient ein \EE{Blutdruckprotokoll}.
  \item[\SAMPLER{L}] Wann war die letzte
  Medikamenteneinnahme bzw. die letzten ärztlichen
  Kontrollen? 
  \item[\SAMPLER{E}] Gibt es Ereignisse die zu einer
  fehlerhaften Einnahme der Medikamente geführt haben
  können?
%   \item[\SAMPLER{R}] 
\end{SAMPLERText}
}
{
\begin{SAMPLERSlide}
  \item[\SAMPLER{S}] Siehe ABCDE.
  \item[\SAMPLER{M}] Blutdrucksenkende Medikamente. 
  Neue Verschreibungen? Einnahmefehler? 
  \item[\SAMPLER{P}] (Chronische) arterielle Hypertonie, Blutdruckprotokoll
  \item[\SAMPLER{L}] Letzte
  Medikamenteneinnahme, letzte ärztliche
  Kontrollen? 
  \item[\SAMPLER{E}] Ursachen für Einnahmefehler?
%   \item[\SAMPLER{R}] 
\end{SAMPLERSlide}
}
{}
{}
{}



% M %%%%%%%%%%%%%%%%%%%%%%%%%%%%%%%%%%%%%%%%%%%%%%%%%%% M %
% Hypertensive Krise
\ProcedurePrintDefault{MI10911C}


% M %%%%%%%%%%%%%%%%%%%%%%%%%%%%%%%%%%%%%%%%%%%%%%%%%%% M %
% Hypertensiver Notfall
\ProcedurePrintDefault{MI10912C}





\clearpage
\section{Störungen des \HeadingKeyword{Herzens}}
\subsection{\HeadingKeyword{Herzinsuffizienz}}
\label{topic:herzinsuffizienz}
\index{Herzinsuffizienz}
\index{Herzinsuffizienz!kompensierte}
\index{Herzinsuffizienz!dekompensierte}
\index{Rechtsherzinsuffizienz}
\index{Linksherzinsuffizienz}
\index{Globalinsuffizienz}
\index{Insuffizienz!Herz-}


\DefinitionInPara{Eine \T{Herzinsuffizienz} bezeichnet eine 
unzureichende Auswurfleistung des Herzens, unabhängig von deren Ursache.}
Es sind viele Ursachen
möglich, manche Ursachen treten \E{akut} auf, manche sind
eher \E{chronischer} Natur.



\TextSlide{Einteilung}
{Die Einteilung der Herzinsuffizienz kann nach
dem betroffenen Herzteil oder der Fähigkeit des Körpers zur
Kompensation erfolgen:
\begin{itemize}
  \item Einteilung nach betroffenen \EE{Herzteil}:
  \begin{itemize}
    \item \T{Rechtsherzinsuffizienz}: Unzureichende Leistung des \E{rechten}
    Herzens. Es kommt zu einem \E{Rückstau} von Blut \E{in
    den Körperkreislauf} und in Folge zu
    \E{Beinödemen} oder
    gestauten Halsvenen.
    \item \T{Linksherzinsuffizienz}:
    Unzureichende Leistung des \E{linken} Herzens. Es kommt
    es zu einem \E{Rückstau} des Blutes \E{in den
    Lungenkreislauf}. Aufgrund der Stauung kommt es zum
    Austritt von Flüssigkeit in das Zwischengewebe und es
    kann ein lebensgefährliches \E{kardiales Lungenödem}
    entstehen.
    Ein typisches Zeichen dafür ist ein
    \I[Atemgeräusch!brodelndes -
    bei Herzinsuffizienz]{brodelndes Atemgeräusch}.
    \item \T{Globalinsuffizienz}: Unzureichende Leistung beider Herzteile.
  \end{itemize}
  \item Einteilung nach \EE{Kompensation}: Abhängig davon, wie
  gut der Körper mit der
  Einschränkung zurecht kommt, unterteilt man:
  \begin{itemize}
    \item \T{Kompensierte Herzinsuffizienz}: Der Körper hat noch
    genug Reserven, kann gegensteuern (kompensieren) und
    kommt mit der eingeschränkten Leistung aus. Für viele
    Patienten ist sie ein \E{chronischer Zustand} mit dem sie
    gut leben können, da sich der Körper mit der Zeit an die
    abnehmende Leistungsfähigkeit gewöhnt.
    \item \T[Dekompensierte
    Herzinsuffizienz]{\E{De}kompensierte Herzinsuffizienz}:
    Problematisch wird es jedoch, wenn die
    \E{Kompensationsmechanismen} (köpereigene, Medikamente,
    {\dots}) versagen und der Patient plötzlich Symptome
    auftreten.
    Der Körper kommt aufgrund des Fortschreitens der Erkrankung 
    oder aufgrund eines erhöhten Bedarfs nicht mehr
    mit der eigenen Leistungsfähigkeit aus.
  \end{itemize}
\end{itemize}


}{
\begin{compactitem}
  \item Herzteil:
  \begin{compactitem}
    \item \E{Rechts}herzinsuffizienz
    \item \E{Links}herzinsuffizienz
    \item \E{Global}insuffizienz
  \end{compactitem}
  \item Kompensation:
  \begin{compactitem}
    \item Kompensierte Herzinsuffizienz
    \item \Ca{De}kompensierte Herzinsuffizienz
  \end{compactitem}
\end{compactitem}
}{}{}{}










\TextSlide
{\TxABCDE}
{~
\label{topic:linksherzinsuffizienz}%
\index{Herzinsuffizienz!Links-}%
\index{Insuffizienz!Linksherz-}%
\index{Linksherzinsuffizienz}%
\index{Lungenödem!kardiales - bei Herzinsuffizienz}%%
\label{topic:rechtsherzinsuffizienz}%
\index{Rechtsherzinsuffizienz}%
\index{Herzinsuffizienz!Rechts-}%
\index{Insuffizienz!Rechtsherz-}%
\begin{ABCDEText}
  %   \item[\ABCDE{1}]
  %%%
  %   \item[\ABCDE{2}]
  %%%
  \item[\ABCDE{3}] Evtl. ist der Patient bleich bzw.
  blass, vielleicht auch kaltschweissig. Mitunter wirken die
  Patienten unruhig und ängstlich.
  %%%
  \item[\ABCDE{4}] Der Patient klagt über {\RedFlag}\,Atemnot (\EE{Dyspnoe}). 
  Oft berichten
  die Patienten, dass sie nur mehr im Sitzen Luft bekommen,
  sie können nicht flach liegen
  (\Z{Orthopnoe}\ZFootnote{\E{Orthopnoe}: Der Patient muss
  aufrecht sitzen um die Atemnot zu lindern.}).
  %%%
  %   \item[\ABCDE{A}]
  %%%
  \item[\ABCDE{B}] Die Atemfrequenz ist
  oft erhöht (\EE{Tachypnoe}\index{Hyperventilation!bei
  Herzinsuffizienz}), oft besteht ein \E{Hustenreiz}. Die
  {\RedFlag}\,Sauerstoffsättigung kann deutlich erniedrigt sein.
  
  Die Linksherzinsuffizienz entsteht aufgrund mangelnder
  Auswurfleistung des linken Herzens. Durch den \E{Rückstau
  von Blut in den Lungenkreislauf} kann es zu einem Lungenödem bis hin zum
  {\RedFlag}\,\EE{schweren kardialen Lungenödem} kommen. 
  %%%
  \item[\ABCDE{C}] Die Herzfrequenz ist
  oft erhöht (\EE{Tachykardie}). 
  Grundsätzlich ist der Blutdruck vermindert (\EE{Hypotonie}).
  Allerdings kann ein erhöhter Blutdruck bei einer
  Blutdruckkrise (\EE{Hypertonie}, \E{hypertensive Krise}
  (\prettyref{topic:hypertensive-krise}), \E{hypertensiver
  Notfall} (\prettyref{topic:hypertensiver-notfall})) der
  eigentliche \E{Auslöser} einer Herzinsuffizenz sein. Der
  Blutdruck
  kann dann natürlich erhöht sein.
  Es können Zeichen eines {\RedFlag}\,kardiogenen
  Schocks auftreten
  (\FullPrettyRef{topic:schock-kardiogener}).
  %%%
%   \item[\ABCDE{D}] 
%   \item[\ABCDE{E}] %
  %%%
  \item[\ABCDE{\ldots}]
  Bei der Rechtsherzinsuffizienz kommt
  es zu einem Rückstau des Blutes in den Körperkreislauf und
  damit zu Ödemen (Beine, \ldots) und gestauten Halsvenen.
  Die Patienten müssen häufig nachts urinieren. \Z{Über längere
  Zeit kann sich die Leber vergrößern und ein sog.
  \enquote{Wasserbauch} (\Z{Aszites}) bilden, dabei wird durch
  den Rückstau Wasser in den freien Bauchraum abgepresst.}
  %%%
  \item[\ABCDE{=}] Es ist wichtig, den Schweregrad und die Kompensation einer
  Herzinsuffizienz zu beurteilen. Man muss die Symptome generell \E{anhand
  ihres Verlaufs beurteilen}: Bekannte, seit längerer Zeit
  bestehende Symptome sind weniger alarmierend als plötzlich
  auftretende Symptome. 
  
  Kritisch sind plötzlich auftretende Symptome, inbesonders ist hier die
  {\RedFlag}\,\EE{Atemnot} zu nennen, weiters eine
  {\RedFlag}\,\EE{niedrige Sauerstoffsättigung} oder Zeichen
  eines {\RedFlag}\,\EE{kardiogenen Schocks}.
  
  Eine besonders kritische
  Komplikation ist das {\RedFlag}\,kardiale \EE{Lungenödem}.
  Wenn es sehr stark ausgeprägt ist, kann man es mit freiem
  Ohr hören (\E{\enquote{brodelndes Atemgeräusch}}). Diese
  Patienten sind als besonders kritisch einzuschätzen, sie
  sind \E{lebensbedrohlich} krank. Das Herz verbraucht oft die
  allerletzten Reserven, schon die Aufregung durch den
  Transport im Stiegenhaus zum Auto kann das Herz endgültig in
  den Herz-Kreislaufstillstand bringen. Es kommt immer wieder
  vor, dass solche Patienten im Stiegenhaus oder im Auto
  reanimationspflichtig werden. Siehe 
  \FullPrettyRef{topic:lungenoedem}.
\end{ABCDEText}
}
{

\begin{ABCDESlide}
  %   \item[\ABCDE{1}]
  %%%
  %   \item[\ABCDE{2}]
  %%%
  \item[\ABCDE{3}] Blass, evtl. kaltschweißig. Evtl.
  agitiert
  %%%
  \item[\ABCDE{4}] \EE{Dyspnoe}, besonders im Liegen
  $\rightarrow$ Sitzende Position!
  
  Evtl. {\RedFlag}\,schweres kardiales Lungenödem 
  %%%
  %   \item[\ABCDE{A}]
  %%%
  \item[\ABCDE{B}] \EE{Tachypnoe}\index{Hyperventilation!bei Herzinsuffizienz}, 
  {\RedFlag}\,Dyspnoe
  %%%
  \item[\ABCDE{C}] \EE{Tachykardie}, \EE{Hypotonie} oder \EE{Hypertonie}
  (\Abk{RR}\,\Sign{↑}\Sign{↑} kann Auslöser sein). Evtl. Zeichen eines {\RedFlag}\,kardiogenen Schocks
  %%%
%   \item[\ABCDE{D}] 
  %%%
%   \item[\ABCDE{E}]
  %%%
  \item[\ABCDE{\ldots}] 
  \begin{itemize}
    \item Lungenödem
    \item {Bein-}/\EE{Hautödeme}
    \item {gestaute Halsvenen}
    \item Nächtliches Wasserlassen
  \end{itemize}
  %%%
  \item[\ABCDE{=}]
  \begin{compactitem}
    \item Chronisch: oft Gewöhnungseffekte, Kompensation
    \item Dekompensation: Plötzlich neu auftretende
    Symptome, evtl. Alarmsymptome
  \end{compactitem}
\end{ABCDESlide}


}
{}
{}
{}





\TextSlide
{{\TxSAMPLER}}
{
\begin{SAMPLERText}
  \item[\SAMPLER{S}] Es werden meist viele Medikamente genommen,
  sowohl für die Herzinsuffizienz, als auch für die
  Begleiterkrankungen.
  \item[\SAMPLER{P}] Die Herzinsuffizienz ist als chronische
  Erkrankung bekannt, entsprechende Befunde vorhanden. Oft
  bestehen Erkrankungen oder gab es Ereignisse, welche
  den Herzmuskel geschädigt haben, wie
  z.\,B. diverse Herzerkrankungen, frühere Herzinfarkte,
  Bluthochdruck, u.\,a.
  \item[\SAMPLER{R}] Auslöser für akute Beschwerden können alle
  Umstände sein, welche vom Herz eine ungewohnte Mehrarbeit
  erfordern: Darunter fallen z.\,B. Infektionen und andere
  Erkankungen, Stress und ein plötzlicher Blutdruckanstieg.
\end{SAMPLERText}
}
{
\begin{SAMPLERSlide}
  \item[\SAMPLER{S}] Viele
  \item[\SAMPLER{P}]  Herzinsuffizienz bereits bekannt,
  \\ Befunde, Bluthochdruck, Herzerkrankungen, frühere
  Herzinfarkte
  \item[\SAMPLER{R}] Ungewohnte Belastung: Infekte,
  Erkrankungen, Blutdruckanstieg
\end{SAMPLERSlide}
}
{}
{}
{}







% M %%%%%%%%%%%%%%%%%%%%%%%%%%%%%%%%%%%%%%%%%%%%%%%%%%% M %
% Herzinsuffizienz, symptomatisch
\ProcedurePrintDefault{MI50090C}






\clearpage % TEMP temp clearpage

\subsection{Erkrankungen der \HeadingKeyword{Herzkrankzgefäße}}
\index{Koronare Herzkrankheit}
\index{Herzkrankheit!Koronare}
\index{KHK}
\index{Akutes Koronarsyndrom}
\index{Koronarsyndrom!Akutes}
\label{topic:khk}
\label{topic:aks}
\label{topic:acs}
\label{topic:akutes-koronarsyndrom}

\TextSlide{\TxBeschreibung}
{\DefinitionInPara{Die \T{Koronare Herzkrankheit} (\T{KHK}) ist eine
Erkrankung der \EE{Herzkranzgefäße} (Koronargefäße) bei der
es zu einer Verengung der Gefäße kommt.} Sie ist eine
\EE{chronische} Erkrankung, welche in akuten Notfallsituationen
resultieren kann. 
\DefinitionInPara{Das 
\T[Akutes Koronarsyndrom]{Akute Koronarsyndrom}
\index{Koronarsyndrom!akutes} (\I{ACS}\ZFootnote{\I{ACS}: \I{Acute Coronary Syndrome}})
ist ein  Symptomenkomplex, welcher
typischerweise bei Erkrankungen wie dem 
\E{Herzinfarkt} und der \E{stabilen} oder \E{instabilen Angina
pectoris} \EE{akut} auftritt.} 
Bei diesen Krankheitsbildern kommt es zu einer akuten \EE{Unterversorgung des
Herzmuskels} mit Blut und Sauerstoff aufgrund der
Durchblutungsstörung. Die Ursache sind in erster Linie
verengte oder verstopfte Herzkranzgefäße
(\I[Ischämie!Akutes Koronarsyndrom]{Ischämie},
\FullPrettyRef{topic:ischaemie}). Das akute Koronarsyndrom
ist häufig das Ergebnis einer vorbestehenden koronaren Herzkrankheit.
}
{%
\begin{itemize}
  \item Erkrankung der Herzkranzgefäße $\rightarrow$ Verengung
  \item Koronare Herzkrankheit (KHK, chronisch)
  \item Akutes Koronarsyndrom (ACS, akut)
  \begin{itemize}
    \item Symptomenkomplex, z.\,B. bei:
    \begin{itemize}
      \item Angina pectoris
      \item Herzinfarkt
    \end{itemize}
  \end{itemize}
\end{itemize}
}
{}
{}
{}
\Commentary[Akutes Koronarsyndrom]{
Das
\protect\T[Akutes Koronarsyndrom]{Akute Koronarsyndrom}
\protect\index{Koronarsyndrom!akutes}
ist ein Symptomenkomplex, welcher
typischerweise bei Erkrankungen wie dem
\E{Herzinfarkt} und der \E{stabilen} oder \E{instabilen Angina
pectoris} auftritt. Daneben gibt es noch andere
Erkrankungen, welche sich von den vorgenannten entweder
durch die Pathogenese (z.\,B. \protect\I{Prinzmetal-Angina}) oder
durch den Pathomechanismus (z.\,B. \protect\I{Tako-Tsubo-Kardiopathie})
wesentlich unterscheiden.
}
\Commentary[Koronare Herzkrankheit]{Eine koronare
Herzkrankheit ist oft anamnestisch bereits bekannt. Ein
Hinweis kann auch die Einnahme von ASS-Präparaten (z.\,B.
ThromboASS{\texttrademark}, HerzschutzASS{\texttrademark},
\ldots) oder ein Stent-Ausweis sein. 
Je nachdem welche Koronargefäße betroffen sind, spricht man
auch von einer 1-, 2- oder 3-Gefäßerkrankung (\Lang{engl.}
1-, 2-, 3-Vessel-Disease; bzw. \Lang{Abkz.} 1-, 2.-, 3.-VD).}

\begin{PictureSeriesWide}{}{Bilderserie:
\EEE{Koronargefäße}}
\PictureSeriesRowWide{3}
{./images/hirtler-aass/Herzvorderflaeche.pdf}
{Das Herz mit seinen Koronargefäßen \SourcePicture{Hirtler}}
{./images/wikipedia-cc-by/Hk_coronary_bionerd.jpg}
{Darstellung der
Herzkranzgefäße während einer Herzkatheteruntersuchung.
Dabei wird über die Leistenarterie ein Katheter
eingebracht und bis knapp vor das Herz zu den Abgängen der
Koronargefäße aus der Aorta vorgeschoben. Anschließend wird ein
Kontrastmittel appliziert um die Gefäße mittels Röntgen sichtbar zu machen.
\SourcePicture{WmCo
\enquote{Bionerd}, CC-BY-3.0}}
{./images/wikipedia-cc-by/Heart_coronary_artery_lesion-lq.jpg}
{Die Koronargefäße versorgen den
Herzmuskel von außen nach innen. \SourcePicture{Patrick J.
Lynch, CC-BY-2.5}} {./images/wikipedia-cc-by/Hk_coronary_bionerd.jpg}
{empty}
{}

\end{PictureSeriesWide}

\TextSlide
{Akutes Koronarsyndrom}
{%
Da bei Behandlungsbeginn oft nicht klar ist, ob es sich um
einen Herzinfarkt oder einen Angina-pectoris-Anfall handelt,
wird bei Vorliegen der entsprechenden Symptome (Atemnot,
Thoraxschmerz, Engegefühl im Thorax, 
Todesangst, \ldots) einfach vom \T{Akuten
Koronarsyndrom}\ZFootnote{\Lang{Engl} Acute Coronary
Syndrome, \Abk{ACS}.} gesprochen. 

Es sind \E{zwei Auslösemechanismen} möglich:
Einerseits kann eine Minderdurchblutung des Herzmuskels
(Myokard)
  aufgrund einer oder mehrerer \E{Einengung(en) von
  Herzkranzgefäßen} (Koronararterien) die Ursache sein, oder
es kommt zu einem \E{erhöhtem Sauerstoffbedarf} des Herzens
bei
  Anstrengung.
In beiden Fällen kommt es zu einem \EE{Missverhältnis
zwischen Sauerstoffangebot und Sauerstoffbedarf} des
Herzens (\I[Ischämie!Angina pectoris]{Ischämie},
\FullPrettyRef{topic:ischaemie}), die Folge sind Schmerzen,
Atemnot und ein Leistungsverlust!
}
{%
\begin{itemize}
  \item Unklare Diagnose, typische Symptome:
\end{itemize}
\begin{itemize}
  \item Atemnot
  \item Thoraxschmerz
  \item Engegefühl im Thorax
  \item \ldots
\end{itemize}
}

\Layout{27}
{Akutes Koronarsyndrom}
{}
{}
{./images/hirtler-aass/Herzschmerz.pdf}
{Typische Schmerzausdehnung bzw. -ausstrahlung beim ischämischen
Herzschmerz.}
{Hirtler.}




\TextSlide
{\TxABCDE}
{~
\begin{ABCDEText}
%   \item[\ABCDE{1}] 
%   %%%
%   \item[\ABCDE{2}] 
  %%%
  \item[\ABCDE{3}]  Der Patient macht oft einen leidenden
  Eindruck, eventuell ist er {\RedFlag}\,kaltschweißig und blass.
  %%%
  \item[\ABCDE{4}] Die Leitsymptome sind ein {\RedFlag}\,\EE{akuter,
  nicht atemabhängiger, Brustschmerz} und {\RedFlag}\,\EE{Atemnot}
  (\E{Dyspnoe}).
  Der Schmerz ist meist hinter dem Brustbein lokalisiert und
  kann \EE{ausstrahlen}, häufig in den Kiefer, (linken) Arm
  und in den Rücken. Der Schmerz wird als brennend bzw.
  stechend beschrieben. Zusätzlich zu dem Schmerz verspürt
  der Patient auch oft ein \EE{Enge-} bzw. \EE{Druckgefühl}
  im Brustkorb, \Speech{\enquote{als würde jemand darauf
  stehen}}.
  Unter Umständen kann es auch zu einem
  \EE{Oberbauchschmerz} kommen, welcher als
  Verdauungsstörung oder Magen-Darm-Erkrankung fehlgedeutet
  werden kann.
  
  Bei Diabetikern kann es aufgrund der Nervenschädigungen
  auch zu einem schmerzlosen Verlauf kommen, man spricht
  dabei auch vom \enquote*{\T[stummer Infarkt]{stummen
  Infarkt}}\index{Infarkt, stummer}.
%   Sonst steht normalerweise der \E{stechende/brennende},
%   nicht atemabhängige \II[Brustschmerz]{Brust-} oder
%   \II{Oberbauchschmerz} zusammen mit \II{Atemnot} und
%   \II{Todesangst} im Vordergrund. Der Brustschmerz kann in
%   Richtung \E{Magen}, rechter oder linker \E{Hand},
%   \E{Kiefer} oder auch \E{Rücken} \EE{ausstrahlen}.
  
  \EE{Todesangst und Vernichtungsgefühl} sind
  häufig.
  Bei Auftreten eines \EE{kardiogenen Schocks} können die
  entsprechenden Schockzeichen beobachtet werden: Blässe,
  Kaltschweißigkeit, Hypotonie, etc.
  (\FullPrettyRef{topic:schock-kardiogener}).
  %%%
  \item[\ABCDE{B}] {\RedFlag}\,\EE{Atemnot} kommt häufig vor und ist ein
  Leitsymptom.  Begleitend kann man oft eine Zyanose
  beobachten.
  %%%
  \item[\ABCDE{C}] Der Puls kann unregelmäßig sein.   
  Der Blutdruck kann in jede Richtung
  verändert sein. Bedenke: Ein hypertensiver Notfall kann
  Auslöser eines \Abk{ACS} sein. Es
  können Zeichen eines {\RedFlag}\,\E{kardiogenen Schocks} vorliegen. 
  
  %%%
%   \item[\ABCDE{D}]
  %%%
  \item[\ABCDE{E}]  Das EKG kann wichtige Hinweise über
  Ausdehnung, Alter und Lokalisation des Infarkts bringen.
  Um die Dynamik beurteilen zu können, soll schon möglichst
  früh ein 12-Kanal-EKG geschrieben werden um Veränderungen
  im Verlauf zu dokumentieren.
  Die Beurteilung ist jedoch nur entsprechend ausgebildeten
  und geübten Personal vorbehalten.
%   %%%
%   \item[\ABCDE{\ldots}] 
  %%%
  \item[\ABCDE{=}]  Ein Patient mit einem  {\RedFlag}\,akuten
  Koronarsyndrom ist grundsätzlich als vital bedroht
  einzuschätzen. Eine ärztliche Diagnose ist auch
  hinsichtlich der Auswahl des richtigen Transportziel
  (Herzkatheterlabor?) entscheidend.
\end{ABCDEText}
}
{
\begin{ABCDESlide}
%   \item[\ABCDE{1}] 
%   %%%
%   \item[\ABCDE{2}] 
  %%%
  \item[\ABCDE{3}]  Leidend, 
  evtl. kaltschweißig, blass
  %%%
  \item[\ABCDE{4}] {\RedFlag}\,\EE{Akuter,
  nicht atemabhängiger, Brustschmerz} (brennend, stechend;
  Hinter dem Brustbein, evtl.
   \EE{Ausstrahlung} in Kiefer, (linken) Arm, in Rücken)
  und
  {\RedFlag}\,\EE{Atemnot}, oft \EE{Enge-} bzw.
  \EE{Druckgefühl}, Todesangst
  
 

  Typisch untypisch:
  \EE{Oberbauchschmerz};  
  \enquote*{{stummer
  Infarkt}}  
  %%%
  \item[\ABCDE{B}] {\RedFlag}\,\EE{Atemnot}, evtl. {\RedFlag}\,Zyanose
  %%%
  \item[\ABCDE{C}] Puls evtl. unregelmäßig,   
  Blutdruck in jede Richtung
  verändert, evtl. Zeichen eines {\RedFlag}\,{kardiogenen Schocks}
  
  %%%
%   \item[\ABCDE{D}]
  %%%
  \item[\ABCDE{E}] EKG (im Verlauf)! 
%   %%%
%   \item[\ABCDE{\ldots}] 
  %%%
  \item[\ABCDE{=}]  Vital bedroht
\end{ABCDESlide}
}
{}
{}
{}
 
\TextSlide
{{\TxSAMPLER} beim akuten Koronarsyndrom}
{
\begin{SAMPLERText}
  \item[\SAMPLER{S}] Siehe {\TxABCDE}
  %%%  
  \item[\SAMPLER{M}] \II{Nitroglycerin} wird zur Erweiterung der
  Koronargefäße von KHK-Patienten eingenommen. Dies kann
  entweder mittels eines Pflasters (Nitro-Pflaster), über
  Kapseln oder mittels eines Sprays (z.\,B.
  \I{Nitrolingual{\texttrademark}}), welcher bei Bedarf
  eingenommen wird, erfolgen. 
  Herzkranke Patienten nehmen daneben oft noch viele andere Medikamente ein.
  \item[\SAMPLER{P}]
  Bei den meisten Patienten ist die Angina
  Pectoris schon bekannt, bzw. es besteht eine \EE{koronare
  Herz-Krankheit} (\E{KHK}). Die KHK ist eine Erkrankung der
  Herzkranzgefäße, diese werden durch Ablagerung von Fetten
  und Kalk geschädigt und eingeengt, dadurch kann weniger Blut
  durchfließen. So kommt es zu einer Mangelversorgung des
  Herzmuskels und zu den Symptomen der Angina pectoris.
  \item[\SAMPLER{E}]
  Häufig wird über körperliche Anstrengung oder Aufregung
  vor Beginn des Anfalls berichtet.
\end{SAMPLERText}
}
{
\begin{SAMPLERSlide}
  \item[\SAMPLER{S}] Siehe {\TxABCDE}
%  \begin{itemize}
%    \item Stechender/brennender \EE{Brustschmerz} mit
%    Ausstrahlung oder \EE{Oberbauchschmerz}
%    
%    \enquote*{Stummer Verlauf} bei Diabetikern möglich
%    \item \EE{Dyspnoe}
%    \item Druckgefühl, \Speech{\enquote{als würde jemand darauf
%    stehen}}
%    \item Todesangst
%    \item Ggfs. kardiogener Schock
%  \end{itemize}
  \item[\SAMPLER{M}] Nitroglycerin: Pflaster, Kapseln, Spray
  \item[\SAMPLER{P}] Koronare Herzkrankheit (KHK)
  \item[\SAMPLER{E}] Körperliche Anstrengung
\end{SAMPLERSlide}
}

\TextSlide
{Angina pectoris}
{\index{Angina pectoris}
\label{topic:angina-pectoris}
\HeadingSub{\Lang{lat.} \enquote{Brustenge}}

\noindent\DefinitionInPara{Plötzliches Engegefühl in der Brust meist
mit Thoraxschmerz infolge einer \E{vorübergehenden}
Unter\-ver\-sorgung des Herzmuskels mit sauerstoffreichem Blut.}
Ein akuter Angina pectoris-Anfall ist primär \EE{nicht
sicher von einem Herzinfarkt unterscheidbar}. Die Symptome
treten eher bei Belastung auf und vergehen oft bei
Entspannung, das Vorübergehen der Beschwerden darf aber
natürlich nicht abgewartet werden!



Der akute Angina-Pectoris-Anfall gehört zu der Familie des
\EE{Akuten Koronarsyndroms}. Dazu gehört auch der
Herzinfarkt, dem der gleiche Mechanismus wie der Angina Pectoris zu
Grunde liegt, wobei es beim Infarkt im \E{Gegensatz} zur Angina
pectoris zu einem \E{totalen} Herzkranzgefäß\E{verschluss}
mit Absterben von Herzmuskelgewebe kommt.
\A{KOCH: nein; GABS: Doch. }}
{%
Missverhältnis zwischen Sauerstoffangebot und  
Sauerstoffbedarf
\begin{itemize}
  \item Einengung(en) von Herzkranzgefäßen
  \item Erhöhter Sauerstoffbedarf bei Anstrengung
  \item Symptome des akuten Koronarsyndrom,  eher bei
  Belastung, vergehen oft bei Entspannung
  \item Primär \EE{nicht sicher vom Herzinfarkt unterscheidbar}
\end{itemize}
}
{}
{}
{}




\TextSlide
{Herzinfarkt}
{\HeadingSub{\Lang{Term.} \T{Myokardinfarkt}, \Lang{Abkz.} \T{MCI}.}
\index{Herzinfarkt}
\index{Infarkt!Herz-}
\index{Myokardinfarkt}
\index{MCI}
\index{Herzkranzgefäße!Verschluss}
\index{Koronargefäße!Verschluss}
\label{topic:herzinfarkt}
\label{topic:myocardinfarkt}
\label{topic:mci}

\noindent\DefinitionInPara{Verschluss eines oder mehrerer Herzkranzgefäße mit
  anschließendem Absterben des betroffenen Herzmuskelgewebes
  (\I[Ischämie!Herzinfarkt]{Ischämie},
\FullPrettyRef{topic:ischaemie}).}
Die Symptome entsprechen denen des akuten Koronarsyndroms, 
doch vergehen diese nicht mehr und es bleiben --
selbst bei optimaler Behandlung -- Schäden am Herzmuskel
bestehen. 


}
{%
\begin{itemize}
  \item \EE{Verschluss} von Herzkranzgefäßen
  \item \EE{Absterben} von Herzmuskelgewebe
  \item Symptome eines \EE{ACS}
\end{itemize}}
{}
{}
{}


\begin{PictureSeriesWide}{}{Bilderserie:
\EEE{Schädigung des Herzmuskel beim Herzinfarkt}}
\PictureSeriesRowWide{3}
{./images/wikipedia-cc-by-sa/Heart_normal_short_axis_section-00640}
{Das Herz im Querschnitt, gesehen von unten: Links der
rechte Ventrikel, rechts der muskelstarke linke Ventrikel.
 \SourcePicture{Patrick Lynch, CC-BY}}
{./images/wikipedia-cc-by-sa/Heart_inferior_wall_infarct-00640}
{Ischämisches Herzmuskelgewebe {\ldots}
\SourcePicture{Patrick J. Lynch, medical illustrator; C. Carl Jaffe, MD, cardiologist; CC-BY 2.5}} 
{./images/wikipedia-cc-by-sa/Heart_inferior_wall_scar}
{{\ldots} stirbt nach einger Zeit ab, es bildet sich eine
(funktionslose) Narbe, die Herzleistung ist in Folge beeinträchtigt. 
\SourcePicture{Patrick J. Lynch, medical illustrator; C. Carl Jaffe, MD, cardiologist; CC-BY 2.5}}
{empty}
{}
\end{PictureSeriesWide}

  
\TextSlide{\TxKomplikationen}
{%
\begin{itemize}
  \item \EE{Herzrhythmusstörungen} (alle Arten), bis hin
  zum Kreislaufstillstand und zur Reanimationspflichtigkeit.
  Besonders häufig kommt es als Frühkomplikation zu plötzlichen, 
  lebensgefährlichen Rhythmusstörungen wie der ventrikulären 
  Tachykardie oder dem Kammerflimmern.
  \item Kardiogener Schock \index{Kardiogener
  Schock!beim Herzinfarkt}\index{Schock!Kardiogener!beim
  Herzinfarkt}
\end{itemize}
}{
\begin{itemize}
  \item \EE{Lebensgefahr!}
  \item Jede Herzrhythmusstörung als Komplikation möglich!
  \item Kammerflimmern als Frühkomplikation häufig!
\end{itemize}

}{}{}{}



\Layout{27}
{}
{}
{}
{./images/wikipedia-pd/12_lead_generated_ventricular_tachycardia}
{Lebensbedrohliche Rhythmusstörungen, wie hier die
ventrikuläre Tachykardie, sind häufige Komplikationen eines
Herzinfarktes.}
{WM/PD}





% M %%%%%%%%%%%%%%%%%%%%%%%%%%%%%%%%%%%%%%%%%%%%%%%%%%% M %
% Thoraxschmerzen, jetzt beschwerdefrei
\ProcedurePrintDefault{MI20091C}


% M %%%%%%%%%%%%%%%%%%%%%%%%%%%%%%%%%%%%%%%%%%%%%%%%%%% M %
% Akutes Koronarsyndrom
\ProcedurePrintDefault{MI24091C}




\TextSlide{Weitere Therapie}
{%
Herzinfarktpatienten benötigen oft eine besondere Behandlung
und normalerweise nicht auf eine beliebige Internistische
Station gebracht\citep{Kalla2006}. Der Notarzt muss
i.\,d.\,R. zwischen dem Einleiten einer Lyse-Therapie und
der Einweisung in ein Herzkatheterlabor entscheiden.

\begin{itemize}
  \item Bei der \II{Lyse-Therapie} versucht man mittels eines
  speziellen Medikaments eine Verstopfung aufzulösen.
  \item \index{PTCA}
  \index{Herzkatheterlabor}
  Mit einem \EE{Herzkatheter}
  versucht man die
  verengten oder verstopften Gefäße aufzudehnen. Dazu
  benötigt man ein Spital, welches über ein entsprechend
  ausgestattetes und
  \E{dienstbereites}
  \T{Herzkatheterlabor} verfügt.\opt{REGVIE}{\footnote{In
  Wien gibt es einen \E{Dienstplan} für die Katheterlabore,
  d.h. jedes Katheterlabor hat an bestimmten Tagen der Woche
  rund um die Uhr Bereitschaftsdienst. Eine Voranmeldung und
  vorangehende Absprache zwischen Notarzt und diensthabendem
  Oberarzt der Abteilung ist erforderlich.}} 
  \hfill\Commentary[Herzkatheter]{%
  Bei
  einer Herzkatheteruntersuchung (\I{Koronarangiographie})
  wird über die Leistenarterie ein Katheter eingebracht und
  bis knapp vor das Herz zu den Abgängen der Koronargefäße
  vorgeschoben und anschließend ein
  Röntgen-Kontrastmittel appliziert. Dadurch kann mittels
  Durchleuchtungsgeräten der Blutfluß in den  Koronargefäßen
  beobachtet und beurteilt werden, sowie Verengungen und
  Verschlüsse identifiziert werden (diagnostische
  Herzkatheteruntersuchung).
  Wird eine solche Stelle gefunden, kann diese oft auch gleich mittels des
  Herzkatheters behandelt werden
  (Intervention; \I{PCI} (Percutane Coronare
  Intervention) oder \I{PTCA} (Perkutane Transluminale
  Coronare Angioplastie)). Dabei wird ein Ballon
  in die Engstelle eingebracht, welcher das Gefäß aufdehnt.
  Anschließend wird ein zylindrisches Metallgeflecht, ein
  sog.
  \I{Stent}, eingebracht, welcher als Stütze für das Gefäß
  dient und dieses längerfristig offen halten soll. Vgl.
  \cite{Harrisons18DeKap230,Harrisons18DeKap246}
  }
\end{itemize}

}
{
\begin{itemize}
  \item Lyse
  \item Herzkatheterlabor (PCI, PTCA)%
  \opt{REGVIE}{%
\begin{itemize}
  \item Dienstbereites Katheterlabor: eigener Dienstplan
\end{itemize}%
}  	
\end{itemize}

}{}{}{}


\begin{figure}[H]
\caption{EKG-Veränderungen beim Herzinfarkt:
Unterschiedliche Gefäße verursachen unterschiedliche
Veränderungen im EKG.} 

\subfloat[\Z{Unauffälliger EKG-Befund}
\SourcePicture{WM/PD}] {\includegraphics[width=\linewidth]{./images/wikipedia-pd/12_lead_generated_sinus_rhythm}}

\subfloat[\Z{EKG-Veränderungen bei
einem Vorderwandinfarkt: ST-Hebungen in den
Brustwandableitungen.}
\SourcePicture{WM/PD}] {\includegraphics[width=\linewidth]{./images/wikipedia-pd/12_lead_generated_anterior_MI}}

\subfloat[\Z{EKG-Veränderungen bei einem Hinterwandinfarkt:
Deutliche ST-Hebungen (\enquote*{Katzenbuckel}) in den
Ableitungen \EKG{II}, \EKG{III} und
\EKG{aVF}}
\SourcePicture{WM/PD}] {\includegraphics[width=\linewidth]{./images/wikipedia-pd/12_lead_generated_inferior_MI}}
\end{figure}



\begin{Literary}
\fullcite{youtube:Herzinfarkt}

\fullcite{youtube:Herzinfarkt2}

\fullcite{Erc2010Sec5}

Sowie:
\cite{Harrisons18DeKap243,Harrisons18DeKap244,Harrisons18DeKap245}
\end{Literary}

\subsection{\HeadingKeyword{Herzrhythmusstörungen}}
\label{topic:herzrhythmusstoerungen}
\index{Herzrhythmusstörungen}
\index{Rhythmusstörungen!Herz-}

Herzrhythmusstörungen sind häufig, allerdings auch oft sehr
unterschiedlich in ihrer Bedeutung: Manche sind chronisch
und eher wenig gefährlich, andere können plötzlich auftreten
und sogar akut lebensbedrohlich sein. Rhythmusstörungen
können sowohl selber das Problem darstellen, als auch \E{als
Komplikationen} im Rahmen anderer Krankheitsbilder (z.\,B.
Herzinfarkt, etc.) auftreten.

\TextSlide{Arten von Herzrhythmusstörungen}
{%
Herzrhythmusstörungen können hervorgerufen werden durch eine
\EE{Störung der Frequenz} (zu langsam, zu schnell) oder
einer \EE{Störung der Regelmäßigkeit} (Rhythmus im engeren Sinn). 
Wenn man die
elektrischen Ströme mittels eines EKG darstellt, kann man
zusätzlich auch \EE{Störungen der Form der elektrischen
Herzströme} erkennen und beschreiben. Ein halbautomatischer
Defibrillator (AED) erledigt dies automatisch, er zeichnet
die elektrischen Herzströme auf und beurteilt diese.

Im Rahmen von Herzrhythmusstörungen kann es zu einer
\EE{Entkoppelung der elektrischen und der mechanischen
Herzaktion} kommen, z.\,B. wenn nicht jedem elektrischen
Impuls ein Herzschlag folgt.
}{
\begin{itemize}
  \item Störung der
  \begin{compactitem}
    \item Frequenz
    \item Regelmäßigkeit
    \item Form der elektrischen Herzströme
  \end{compactitem}
\end{itemize}

Entkoppelung von elektrischer und der mechanischer
Herzaktion möglich!
}{}{}{}


\TextSlide{Hier zählt die Geschwindigkeit: Tachy- und Bradykardie}
{%
%
\begin{itemize}
  \item \T{Bradykardie}:  Langsamer Herzschlag
  ({HF \textless} 60\,\nicefrac{}{\Min} {\tiny (beim Erwachsenen)})
  
  Leistungssportler können aufgrund ihres trainierten Körpers auch sehr
  niedrige Ruhepulsfrequenzen haben.
  \item \T{Tachykardie}:  Schneller
  \marginpar{{\footnotesize Tachometer im Auto:
  Geschwindigkeitsanzeiger, meistens zu schnell ;)}}
  Herzschlag
  (\Abk{HF} {\textgreater} 100\,\nicefrac{}{\Min} {\tiny
  (beim Erwachsenen)})
\end{itemize}

Bei der Beurteilung der Herzfrequenz muss auch immer an die Umstände 
(Ruhe, Belastung, \ldots) und an den
Erregungszustand des Patienten gedacht werden!
}{
\begin{itemize}
  \item Bradykardie \dotfill\ zu langsam
  \item Tachykardie \dotfill\ zu schnell
\end{itemize}
}{}{}{}


%%% Einzelne Tabellen
% 
% \Table
% 	{<Titel>}%1
% 	{<Beschreibung>}%2
% 	{<Label>}%3
% 	{<Quelle>}%4
% 	{<Lizenz>}%5
% 	{<Tabelle>}%6
% 	{<Float>}%7
\TableWide
{Wichtige Herzrhythmusstörungen}
{}
{}
{}
{\AASS}
{% Index-Einträge für Tabelle ---
\label{tab:herzrhythmusstoerungen-wichtige}%
\label{tab:herzrhythmusstoerungen}%
\index{Bradykardie}%
\index{Tachykardie}%
\index{Kammerflimmern}%
\index{Flimmern!Kammer-}%
\index{Vorhofflimmern}%
\index{Flimmern!Vorhof-}%
\index{Asystolie}%
\index{Pulslose Ventrikuläre Tachykardie}%
\index{Tachykardie!Pulslose Ventrikuläre}%
% --- Ende
\noindent\begin{tabulary}{\textwidth}{LLL}\addlinespace\toprule\addlinespace
\noindent\TabHeading{Rhythmus}
& 
\noindent\TabHeading{Bedeutung}
& 
\noindent\TabHeading{Bemerkung}
\\\addlinespace\midrule\addlinespace
\noindent\TabSub{Bradykardie}

{\scriptsize (HF\,{\textless}\,60\PerMin*)}
& 
Situations\-abhängig 
& 
\emph{\enquote{zu langsam}}
\\\addlinespace
\noindent\TabSub{Tachykardie} 

{\scriptsize (HF\,{\textgreater}\,100\PerMin*) }
& 
Situations\-abhängig 
& 
\emph{\enquote{zu schnell}}
\\\addlinespace\midrule\addlinespace
\noindent\TabSub{Kammerflimmern}
& Reanimation  & \emph{\enquote{schockbarer}} Rhythmus 
\\\addlinespace
\noindent\TabSub{Pulslose Ventrikuläre Tachy\-kardie} 
& Reanimation 
& \emph{\enquote{schockbarer}} Rhythmus 
\\\addlinespace
\noindent\TabSub{Asystolie} 

{\scriptsize (HF\,=\,0)} 
& 
Reanimation 
& 
Nulllinie, \emph{\enquote{nicht-schockbarer}} Rhythmus
\\\addlinespace
\noindent\TabSub{Vorhofflimmern}
& regellose Erregung im Vorhof

oft symptomlos 
& 
\E{Chronische Erkrankung.} Sehr viele Leute haben
Vorhofflimmern und leben ganz gut damit. 

Risikofaktor für:

\E{Tachykardie-Anfälle}

\E{Thrombenbildung}, dadurch erhöhtes Lungenembolie- und
Schlaganfallrisiko. Deshalb Vorbeugung mit
gerinnungs\-hemmenden (\enquote*{blut\-ver\-dünnenden})
Medi\-ka\-menten
(z.\,B. Marcoumar{\texttrademark}) ${\rightarrow}$ Patient ist
blutungsgefährdet
\\\addlinespace\midrule\addlinespace
\noindent\TabSub{Extrasystole}
&
Extraschläge
&
Vereinzelt bis \enquote{Salven} möglich

Von symptom-/bedeutungslos bis lebensbedrohlich
\\\addlinespace\bottomrule
\end{tabulary}}
{H}





\subsubsection{Besondere Rhythmen}


\TextSlide{Reanimationspflichtige Rhythmusstörungen}
{~
\begin{itemize}
  \item \T{Pulslose Ventrikuläre Tachykardie} (\I{pVT}):
%   \index{Pulslose Ventrikuläre Tachykardie|TLike}
  \label{topic:pulslose-ventrikulaere-tachykardie}
  %
  Vorstufe zum Kammerflimmern. Das Herz kontrahiert
  so schnell, dass es sich nicht mehr füllen kann
  ${\rightarrow}$ \EE{keine Auswurfleistung}.
  \item \T{Kammerflimmern} (\I{AF}):
%   \index{Kammerflimmern|TLike}
  \label{topic:kammerflimmern}
  %
   Die elektrische
  Erregung im Herzen ist ungerichtet, dadurch kontrahiert der Herzmuskel  nicht mehr geordnet,
  er kommt nur noch zu einem Zittern \EE{ohne nennenswerte
  Auswurfleistung}! (Kammerflimmern darf nicht mit Vorhofflimmern verwechselt werden!)
  \item \T{Asystolie}:
  \label{topic:asystolie}
  Im Herz entsteht keine elektrische
  Erregung, der Herzmuskel kontrahiert daher auch nicht und es erfolgt
  {keine Auswurfleistung}
  \item \T{Pulslose elektrische Aktivität} (\I{PEA}): Bei der
  pulslosen elektrischen Aktivität ist die Herzaktion von
  der elektrischen Aktivität \EE{entkoppelt}, d.\,h. das Herz
  reagiert nicht auf die Impulse des Reizleitungssystems. Der
  EKG-Befund kann zwar unauffällig sein, aufgrund der
  fehlenden Herzaktion besteht trotzdem ein
  Kreislauffstillstand.
\end{itemize}
}
{
\begin{compactitem}
  \item Pulslose Ventrikuläre Tachykardie
  \begin{compactitem}
    \item Keine Zeit zum Füllen ${\rightarrow}$ kein Auswurf
  \end{compactitem}
  \item Kammerflimmern
  \begin{compactitem}  
    \item Wirre Erregung, keine sinnvolle Muskelarbeit ${\rightarrow}$ kein
    Auswurf.
  \end{compactitem}
  \item Asystolie
  \begin{compactitem}
    \item Keine elektrische Erregung, keine Muskelarbeit, kein
    Auswurf
  \end{compactitem}
  \item Pulslose elektrische Aktivität
  \begin{compactitem}
    \item Entkoppelung von Herzschlag und Impulsen des Reizleitungssystems
  \end{compactitem}
\end{compactitem}

}{}{}{}

\Layout{27}
{}
{}
{}
{./images/wikipedia-pd/Lead_II_rhythm_generated_sinus_bradycardia}
{Unauffälliger EKG-Befund}
{WM/PD}

\begin{figure}[H]
\caption{Reanimationspflichtige Rhythmen}
\subfloat[\Z{Ventrikuläre Tachykardie} \SourcePicture{WM/PD}]
{\includegraphics[width=\linewidth]{./images/wikipedia-pd/Lead_II_rhythm_ventricular_tachycardia_Vtach_VT}}

\subfloat[\Z{Kammerflimmern}
\SourcePicture{WM/PD}]
{\includegraphics[width=\linewidth]{./images/wikipedia-pd/Lead_II_rhythm_generated_ventricular_fibrilation_VF}}

\subfloat[\Z{Asystolie}
\SourcePicture{WM/PD}]
{\includegraphics[width=\linewidth]{./images/wikipedia-pd/Lead_II_rhythm_generated_asystole}}
\end{figure}



    
\TextSlide{Vorhofflimmern}
{
\index{Vorhofflimmern|TLike}
\label{topic:vorhofflimmern}
Regellose Erregung im \E{Vorhof}, die oft symptomlos ist. 
(Die Herzkammern arbeiten dabei normal!) Vorhofflimmern kann
ständig bestehen oder auch nur manchmal (episodenhaft)
auftreten. Sehr viele Leute haben Vorhofflimmern und leben
ganz gut damit.
Häufige \EE{Komplikationen} sind jedoch:
\begin{compactitem}
  \item \EE{Tachykardie-Anfälle}
  \item \EE{Thrombenbildung} in den Vorhöfen: Dies stellt wiederum ein Risko dar für:
  \begin{compactitem}
    \item Arterielle Gefäßverschlüsse (\FullPrettyRef{topic:arterieller-gefaessverschluss})
    \item Schlaganfälle (\FullPrettyRef{topic:insult})
    \item Mesenterialinfarkte (\FullPrettyRef{topic:mesenterialinfarkt})
  \end{compactitem}
  Daher werden diesen Patienten meist \EE{gerinnungshemmende}
  (\enquote*{blutverdünnenden}) Medikamente zur Vorbeugung verschrieben (z.\,B. \I[Marcoumar!Vorhofflimmern]{Marcoumar}, der Patient ist blutungsgefährdet)
\end{compactitem}
}
{%
\begin{compactitem}
  \item Flimmern im \EE{Vorhof}
  \item Häufige und oft chronische Erkrankung
  \item Thrombenbildung, Tachykardieanfälle
  \item Oft \E{gerinnngshemmende Medikamente}
\end{compactitem}

}{}{}{}


\TextSlide{Extrasystolen}
{%
\index{Extrasystolen|TLike}
\label{topic:extrasystolen}
Extrasystolen sind \enquote{Extraschläge} des Herzens, sie können
vereinzelt oder gehäuft auftreten, eventuell auch
\enquote{salvenartig}. Je nach Häufigkeit haben Extrasystolen
unterschiedliche Auswirkungen auf den Kreislauf und können
symptom- bzw. bedeutungslos sein, oder auch Beschwerden
verursachen. Wenn die Extrasystolen die normale
Herzpumpfunktion massiv stören, können Sie auch
lebensbedrohlich sein.
}
{%
\begin{itemize}
  \item Extraschläge
  \item Vereinzelt od. gehäuft, evtl. \enquote{salvenartig}
  \item Unterschiedliche Auswirkungen auf Kreislauf: symptomlos bis
  lebensbedrohlich möglich
\end{itemize}

}{}{}{}





\subsubsection{\HeadingKeyword{Tachykarde Attacke}}
\label{topic:tachykarde-attacke}


\TextSlide{\TxBeschreibung}
{%
\DefinitionInPara{Plötzliche Tachykardie ohne erkennbare Ursache, wobei der Patient
(meistens) Symptome wahrnimmt und eine Herzfrequenz um bzw. über 160\PerMin*
vorliegt.} 
Tachykarde Attacken sind ein häufiges Notfallbild. Oft sind
chronische Erkrankungen die Ursache, diese
können jedoch sehr unterschiedlich sein und sind auch vom Fachmann bei weiterführenden Untersuchungen
oft nicht leicht zu ermitteln.
}{
\begin{itemize}
  \item HF\,$\geq$\,160\PerMin*
\end{itemize}
}{}{}{}

\TextSlide
{\TxABCDE}
{
\begin{ABCDEText}
%   \item[{\color{AltBlueLight}\usefont{T1}{PTSansNarrow-TLF}{m}{n}\selectfont A-Problem}] 
  \item[\ABCDE{B}] {\RedFlag\,}Dyspnoe kann in Folge der kardialen Mehrbelastung auftreten
  \item[\ABCDE{C}] Es liegt eine \EE{Tachykardie} vor. Der Puls
  ist i.\,d.\,R. hoch, meist {\RedFlag\,}über 160\PerMin* und evtl. nur
  schwach tastbar.
  Der Patient kann bei mangelnder Herzauswurfleistung
  {\RedFlag\,}kreislaufinsuffizient werden und einen {\RedFlag\,}kardiogenen Schock
  entwickeln. Im Extremfall ist auch ein {\RedFlag\,}Kreislaufstillstand
  möglich.
%   \item[{\color{AltBlueLight}\usefont{T1}{PTSansNarrow-TLF}{m}{n}\selectfont D-Problem}]
%   \item[{\color{AltBlueLight}\usefont{T1}{PTSansNarrow-TLF}{m}{n}\selectfont E-Problem}]
\end{ABCDEText}
}
{
\begin{ABCDESlide}
%   \item[{\color{AltBlueLight}\usefont{T1}{PTSansNarrow-TLF}{m}{n}\selectfont A-Problem}] 
  \item[\ABCDE{B}] Evtl. {\RedFlag\,}Dyspnoe
  \item[\ABCDE{C}] Tachykardie {\RedFlag\,}$\geq$\,160\PerMin*, Puls evtl. schwach tastbar; 
  
  Evtl. {\RedFlag\,}Kreislaufinsuffizienz, {\RedFlag\,}kardiogener Schock, {\RedFlag\,}Kreislaufstillstand
%   \item[{\color{AltBlueLight}\usefont{T1}{PTSansNarrow-TLF}{m}{n}\selectfont D-Problem}] 
%   \item[{\color{AltBlueLight}\usefont{T1}{PTSansNarrow-TLF}{m}{n}\selectfont E-Problem}]
  
\end{ABCDESlide}
}
{
}
{}
{}



\TextSlide
{{\TxSAMPLER}}
{
\begin{SAMPLERText}
  \item[\SAMPLER{S}] Das Leitsymptom ist die \EE{Tachykardie} vor. 
  Der Patient
  klagt oft über ein \EE{\enquote{Herzklopfen}} und ist unruhig. Manchmal wird der Anfall von
  Angstgefühlen begleitet.
  
  Oft kommen \EE{Begleitsymptome} dazu, wie zum Beispiel Atemnot oder ein
  leichter Brustschmerz. Dies sind erste Zeichen, dass das Herz an
  die Grenze seiner Belastbarkeit stößt.
  \item[\SAMPLER{M}] Viele Patienten, bei denen
  Rhythmusstörungen bekannt sind nehmen entsprechende
  Medikamente \Z{(Antiarrhythmika)}. \Z{Bei manchen
  Rhythmusstörungen ist es üblich, dem Patienten bestimmte
  Medikamente zu verschreiben, welche er im Anfall nehmen
  soll um diesen zu beenden \Z{(\enquote*{pill in the
  pocket})}.}
  Das Erheben der tatsächlich eingenommenen sowie der verschriebenen Medikamente
  ist sehr wichtig, da diese Informationen des weiteren Behandlungsverlauf
  stark beeinflusst!
  \item[\SAMPLER{P}] Eine große Rolle spielt die \EE{Anamnese}: Da oft
  chronische Erkrankungen die Ursache sind, haben die
  Patienten diese Attacken immer wieder und kennen die
  Symptome bereits (und oft auch deren Behandlung). Oft sind die
  entsprechenden Grunderkrankungen schon seit langem bekannt und es 
  gibt bereits entsprechende \EE{Arztbriefe} oder \EE{Befunde}.
  \item[\SAMPLER{L}] Bei regelmäßig wiederkehrenden Beschwerden sind 
  Informationen über die vorherigen Episoden einzuholen.
  \item[\SAMPLER{E}] Die Ereignisse sind genau zu erfragen und zu
  dokumentieren, eventuell ergibt sich dadurch eine Spur zum auslösenden Faktor.
\end{SAMPLERText}
}
{
\begin{SAMPLERSlide}
  \item[\SAMPLER{S}]
  \begin{itemize}
    \item Tachykardie $\geq$\,160\PerMin*
    \item Spürbares Herzklopfen
    \item Evtl. Atemnot, leichter Brustschmerz, Angstgefühl
  \end{itemize}
  \item[\SAMPLER{M}] \Z{Antiarrhythmika.} Eingenommene und verschriebene Medikamente unbedingt sorgfältig erheben!
  \item[\SAMPLER{P}] Oft vorbekannte Erkrankung(en); Arztbriefe und Befunde!
  \item[\SAMPLER{L}] Frühere Episoden?
  \item[\SAMPLER{E}] Genau erfragen und dokumentieren!
\end{SAMPLERSlide}
}
{}
{}
{}



\TextSlide{{\Lightning}Gefahr }
{%
Da das \EE{Herz plötzlich viel mehr Arbeit leisten muss},
kann es sein, dass es an seine \EE{Grenzen} stößt und
\EE{versagt} bzw. die Blutversorgung des Herzens für diese
Arbeit nicht mehr reicht und ein Angina-pectoris-Anfall oder
sogar ein Herzinfarkt\A{ KOCH: passt mit der Definition MCI
nicht überein, AP wäre besser oder? GABS: AP dazugenommen,
Infarkt aber gut möglich. } ausgelöst wird.
Weiters kann die Rhytmusstörung in eine
andere, akut lebensbedrohende Rhythmusstörungen umschlagen
und es dadurch zu einer akut lebensbedrohlichen Situation
kommen.
\Commentary[Tachykarde Attacke / Gefahren]
{Es gibt viele verschiedene Auslöser für tachykarde
Attacken. Manche Rhythmusstörungen sind tückisch, dass sie
von einem Moment umschlagen und einen Kreislaufstillstand
verursachen können. Zum Beispiel beim Kammerflattern besteht
oft eine \Conc{2}{1}-Überleitung vom Vorhof in die Kammern
mit einer resultierenden Frequenz um die 140\PerMin*. Kommt es
zu der gefürchteten \Conc{1}{1}-Überleitung wird die
Flatterfrequenz der Vorhöfe von \Range{250}{350}\PerMin* (!) direkt auf die
Kammern übergeleitet, dies kommt meist einem Herzstillstand
gleich. \cite{BLIM4}}

\bigskip
}{
\begin{itemize}
  \item Herz überfordert:
  \begin{compactitem}
    \item Angina pectoris
    \item Herzinfarkt
  \end{compactitem}
\end{itemize}
}{}{}{}



\clearpage % TEMP
% M %%%%%%%%%%%%%%%%%%%%%%%%%%%%%%%%%%%%%%%%%%%%%%%%%%% M %
% Tachykarde Attacke
\ProcedurePrintDefault{MR47091C}


\bigskip

\section{Störungen der \HeadingKeyword{Gefäße}}

\subsection{\HeadingKeyword{Aneurysmen} sind Gefäßaussackungen}
\label{topic:aneurysma}

\Text
{{\TxBeschreibung}: Aneurysma allgmein}
{%
\DefinitionInPara{Als \T{Aneurysma} bezeichnet man eine
Aussackung und Erweiterung von Blutgefäßen aufgrund von
Veränderungen der Gefäßwand.} 
Durch die
zunehmende Erweiterung oder Verletzungen der
Gefäßwandschichten kann es zu einem Riss des Aneurysmas und
zu Blutungen kommen. Aneurysmen treten häufig an den großen
Gefäßen (Aorta, Leistenarterien) und an den Hirngefäßen auf.
Folgen eines geplatzen (rupturierten) Aneurysmas können
z.\,B. ein hypovolämischer Schock oder ein blutiger Insult
sein.}
{
\begin{itemize}
  \item Aussackung eines Blutgefäßes
  \item Z.\,B. Aorta, Hirngefäße
  \item Gefahr: Schock, Insult
\end{itemize}
}
{}
{}
{}

\subsubsection{\HeadingKeyword{Aortenaneurysma}}

\label{topic:thoraxaortenaneurysma}
\label{topic:bauchaortenaneurysma}
\nocite{Meron2004}

\TextSlide{\TxBeschreibung{} und Ursachen}
{%%% Gerhard Gabriel <gerhard.gabriel@grissu.org>
Aneurysmen können im gesammten Verlauf der Aorta auftreten, man unterscheidet
zwischen dem \T[Aortenaneurysma!thorakales]{thorakalen} und
den \T[Aortenaneurysma!abdominelles]{abdominellen
Aortenaneurysma} (\I{Bauchaortenaneurysma}).  Zumeist
ist eine krankhafte Veränderung der Aortenwand im
Rahmen von Begleiterkrankungen (Diabetes mellitus, Hypertonie,
Arteriosklerose,
Bindegewebserkrankungen) oder ein Trauma
verantwortlich. Die Aussackung kann alle Gefäßschichten
umfassen oder es können einzelne Schichten der
Gefäßwand einreißen und eine (verletzliche) Ausbuchtung
bilden. 
Die akute Gefahr ist eine
Ruptur. 

Oft sind Aneurysmen asymptomatisch und werden als
Zufallsbefund z.\,B. bei Röntgenuntersuchungen entdeckt. Da
das Operationsrisiko sehr hoch ist werden Aneurysmen erst ab
einer gewissen Größe operiert, bei kleinen Formen werden nur
regelmäßige Kontrollen durchgeführt.

\subparagraph{Anmerkung} Aortenaneursymen sind sehr oft
schwer zu diagnostizieren, in der präklinischen Versorgung
kann im Regelfall nur durch die Anamnese bei einem bereits
bekannten Aneurysma ein entsprechender Verdacht geäussert
werden -- eine sichere Diagnose ist praktisch kaum möglich.
Auch in der klinischen Behandlung werden Aneurysmen häufig
erst sehr spät erkannt, da die Symptome oft eher auf andere
Erkrankungen hinweisen (Lumbago, ACS, Koliken, \ldots). 

}
{%
\begin{itemize}
  \item Gefäßwandaussackung
  \item Thorakales und abdominelles Aortenaneurysma
  \item Evtl. vorbekannt
\end{itemize}
}{}{}{}

% \TextSlide{\TxSymptome}{%
% Je nach
% Größe der Schädigung und entsprechendem Schweregrad der Blutung
% im Läsionsbereich, kommt es binnen kürzester Zeit unter
% heftigem, akutem Bauchschmerz zum Tod. Kleinere Läsionen
% sind ebenfalls von heftigem, akut einsetzendem Bauchschmerz
% begleitet. Es bildet sich ein Blutungsschock (hypovolämischer
% Schock) aus, der unbehandelt zum Tode führt.
% }{
% \begin{itemize}
%   \item Heftigste Schmerzen
%   (\enquote{Vernichtungsschmerz})
%   \item Hoher Blutverlust ${\rightarrow}$ schnell im
%   Schock!
%   \item Evtl. bereits bekanntes Aneurysma
%   ${\rightarrow}$\,Anamnese
%   \item Evtl. tastbare pulsierende Schwellung
% \end{itemize}
% }{}{}{}

\TextSlide
{\TxABCDE}
{~
\begin{ABCDEText}
  \item[\ABCDE{2}] Evtl. können
  {\RedFlag\,}Schockzeichen vorliegen (Blässe, Kaltschweißigkeit,
  \ldots).
  Manche Patienten nehmen eine Schonhaltung aufgrund der
  Schmerzen ein.
  \item[\ABCDE{4}] Je nach Ort des
  Aneurysmas klagt der Patient über {\RedFlag\,}\EE{Brust-}, \EE{Bauch-}
  oder \EE{Rückenschmerzen}, diese werden oft als schneidend
  und plötzlich einsetzend beschrieben und sind oft (aber nicht immer) heftig
  (\I{Vernichtungsschmerz}). Bei thorakalen
  Aortenaneurysmen kommt es häufig zu Symptomen eines
  {\RedFlag\,}\EE{Akuten Koronarsyndroms} (\I[ACS!bei thorakalem
  Aortenaneurysma]{ACS}, \FullPrettyRef{topic:acs}).
  \item[\ABCDE{C}] Je nach Umfang der Blutung kann
  der Patient völlig \EE{kreislaufstabil} sein oder
  {\RedFlag\,}\EE{Schockzeichen} zeigen (hypovolämischer Schock).
  Herzfrequenz und Blutdruck können entsprechend eines
  Schockgeschehens verändert sein.
  \item[\ABCDE{E}] Bei einer Blutung in den Bauchraum 
  kann es zu einer Verhärtung 
  der Bauchdecke kommen ({\RedFlag\,}Akutes Abdomen, 
  \FullPrettyRef{topic:akutes-abdomen}). 
  \Z{Selten kann eine deutliche pulsierende
  Schwellung bei einem abdominellen Aortenaneurysma im Bauch
  getastet werden.}
\end{ABCDEText}
}
{
\begin{ABCDESlide}
  \item[\ABCDE{2}] Evtl. {\RedFlag\,}Schockzeichen, Schohnhaltung
  \item[\ABCDE{4}] {\RedFlag\,}Brust-, Bauch- oder Rückenschmerzen;
  evtl. Symptome eines {\RedFlag\,}ACS
%   \item[\ABCDE{B}] 
  \item[\ABCDE{C}] Kreislaufstabil, evtl. {\RedFlag\,}Schockzeichen
  \item[\ABCDE{E}] {\RedFlag\,}Harte Bauchdecke bei Blutung in Bauch
\end{ABCDESlide}
}
{}
{}
{}

\TextSlide
{{\TxSAMPLER}}
{
\begin{SAMPLERText}
  \item[\SAMPLER{S}] Siehe ABCDE. Evtl. haben die ersten
  Beschwerden vor einigen Tagen oder Wochen z.\,B. in Form
  von leichten Rückenschmerzen begonnen. 
  \item[\SAMPLER{P}] Mitunter ist bereits das Aneurysma
  bekannt (Befunde, Arztbriefe!).
  \item[\SAMPLER{R}] Besonders betroffen sind Patienten mit
  systemischen Kreislauferkrankungen (Hypertonie, erhöhte
  Blutfette, \ldots), sowie Patienten mit einer krankhaften
  Bindegewebsschwäche \Z{(Z.\,B. Marfan-Syndrom)}.
\end{SAMPLERText}
}
{
\begin{SAMPLERSlide}
  \item[\SAMPLER{S}] S. ABCDE; evtl. Rückenschmerzen
  \item[\SAMPLER{P}] Aneurysma evtl. vorbekannt 
  \item[\SAMPLER{R}] Kreislauferkrankungen,
  Bindegewebsschwäche
\end{SAMPLERSlide}
}
{}
{}
{}



\bigskip

% \clearpage
\subsection{\HeadingKeyword{Gefäßverschlüsse} in den Extremitäten}
\label{topic:gefaessverschluesse-extremitaeten}

\TextSlide{Unterscheidung}
{%
Grundsätzlich unterscheidet man je nach betroffenem Gefäß
zwischen \EE{arteriellen} und \EE{venösen}
Gefäßverschlüssen. 
Beim \E{arteriellen Verschluss} besteht
das Hauptproblem in der Unterversorgung der Extremität mit
Blut. Beim \E{venösen Verschluss} kommt es zu einer Ab- und
Rückflussbehinderung. Zusätzlich besteht die Gefahr, dass sich
ein Gerinnsel losreißt und in einem anderen Teil des
Körpers Schaden anrichtet\footnote{Zum Beispiel in der
Lunge im Rahmen einer Lungenembolie}.
}{
\begin{itemize}
  \item \textbf{\textcolor[rgb]{0.7882353,0.0,0.08627451}{Arteriell}}
  \item \textbf{\textcolor{blue}{Venös}}:
  \index{tVT}\T{tVT} (tiefe Venen Thrombose)
\end{itemize}
}{}{}{}



\subsubsection{Peripherer \HeadingKeyword{arterieller Gefäßverschluß}}
\label{topic:arterieller-gefaessverschluss}

\TextSlide{Beschreibung}
{%
Beim \E{arteriellen Gefäßverschluss} wird das
zuführende Gefäß verschlossen. Es kommt zu einer schmerzhaften 
Unterversorgung der Extremität mit Blut
(\I[Ischämie!Arterieller peripherer
Gefäßverschluß]{Ischämie},
\FullPrettyRef{topic:ischaemie}). Der komplette Verschluss passiert plötzlich. }
{
\begin{itemize}
    \item Verschluß eines blut\E{zuführenden} Gefäßes
\end{itemize}
}{}{}{}

\Text{Ursachen}
{%
Meistens liegt eine 
chronische Schädigung der
Gefäße infolge anderer Erkankungen oder aufgrund eines
ungesunden Lebensstils vor. Ursachen für solch eine
Gefäßschädigung sind z.\,B. die Zuckerkrankheit (\T{Diabetes
mellitus}, s.\ref{topic:diabetes}), Hypertonie, erhöhte
Blutfette oder das Rauchen. Die Gefäße verkalken und
verändern sich auch sonst nicht zu ihrem Vorteil.
}
{
\begin{itemize}
  \item Chronische Gefäßschädigung, z.\,B. Bei Diabetes mellitus, Hypertonie \ldots
\end{itemize}
}{}{}{}



\TextSlide{Chronische Grunderkrankung: pAVK}{
Wenn die Gefäße nur verengt sind, spricht man von
der \E{peripheren
arteriellen Verschluss\-krankheit}, abgekürzt \T{paVK}. Sie ist eine
häufige chronische Erkrankung. (\E{Achtung}: Obwohl es
\enquote*{\emph{Verschluss}krankheit} heißt, sind die Gefäße
nicht komplett verschlossen, sondern nur verengt!)
}{
\begin{itemize}
    \item Periphere arterielle Ver\-schluss\-krank\-heit
\end{itemize}

}{}{}{}

 

\TextSlide
{\TxABCDE}
{~
\begin{ABCDEText}
  \item[\ABCDE{2}] Der Patient erscheint
  schmerzgeplagt und nimmt oft eine Schohnhaltung ein.
  \item[\ABCDE{4}] Leitsymptom ist der
  Schmerz, zusammen mit Blässe, Empfindungsstörungen und evtl.
  auch Lähmung der betroffenen Extremität, keine Angabe
  eines Traumas.
  \item[\ABCDE{C}] In der betroffenen Extremität sind die
  peripheren Pulse schlecht oder überhaupt nicht tastbar.
%   \item[\ABCDE{D}] 
%   \item[\ABCDE{E}] 
\end{ABCDEText}
}
{
\begin{ABCDESlide}
%   \item[\ABCDE{1}] 
  \item[\ABCDE{2}] Schmerzen, Schonhaltung
%   \item[\ABCDE{3}] 
  \item[\ABCDE{4}] Lokalisierter Schmerz; Blässe, Empfindungsstörungen, evtl. Lähmung der betr. Extremität
%   \item[\ABCDE{A}] 
%   \item[\ABCDE{B}] 
  \item[\ABCDE{C}] Periphere Pulse der betr. Extremität schlecht/nicht tastbar
%   \item[\ABCDE{D}] 
%   \item[\ABCDE{E}] 
%   \item[\tiny{\fbox{EE}}] test
\end{ABCDESlide}
}
{}
{}
{}

\begin{Synopsis}
  \SynopsisItem{Die klassischen Symptome des arteriellen Gefäßverschlusses sind 
  Schmerz, Blässe,
  Pulslosigkeit,
  Gefühllosigkeit und
  evtl. Lähmung
  der betroffenen Extremität.}
\end{Synopsis}

\TextSlide
{{\TxSAMPLER}}
{
\begin{SAMPLERText}
  \item[\SAMPLER{S}] Siehe \TxABCDE.
  \item[\SAMPLER{M}] Bei Vorhofflimmern wird dem Patienten
  meist eine gerinnungshemmende vorbeugende Therapie
  verordnet (\enquote*{Blutverdünnung}, z.\,B.
  Marcoumar{\texttrademark}, \enquote*{Thrombosespritzen},
  z.\,B. Lovenox{\texttrademark}, Fragmin{\texttrademark}).
  \item[\SAMPLER{P}] Eine periphere arterielle
  Verschlusskrankheit ist häufig schon bekannt. Oft liegen
  Herz-Kreislauf-Erkrankungen oder Diabetes mellitus vor.
  Vorhofflimmern kann aufgrund der Bildung eines Embolus
  einen Verschluss auslösen.
  \item[\SAMPLER{R}] Darunter fallen
  Herz-Kreislauferkrankungen welche die Gefäße schädigen wie
  die arterielle Hypertonie; weiters ist hier der 
  Diabetes mellitus zu nennen.
  Vorhofflimmern kann, wie bereits erwähnt, zur Bildung
  eines Embolus, welcher den Verschluss auslöst, führen,
  insbesonders bei fehlender oder unzureichender
  gerinnungshemmender Therapie.
\end{SAMPLERText}
}
{
\begin{SAMPLERSlide}
  \item[\SAMPLER{S}] \TxABCDE
  \item[\SAMPLER{M}] Je nach Begleiterkrankungen
  \item[\SAMPLER{P}] pAVK, 
  Herz-Kreislauf-Erkrankungen, Diabetes mellitus, 
  Vorhofflimmern 
  \item[\SAMPLER{R}] Siehe \SAMPLER{P}
\end{SAMPLERSlide}
}
{}
{}
{}

\TextSlide{\TxSymptome}
{%
Die Symptome ergeben sich aus der Unterversorgung der
Extremität mit sauerstoffreichem Blut: sie ist blass,
kalt, tut höllisch weh und die Rekap-Zeit ist nicht
mehr sinnvoll messbar. Zusätzlich wird die Extremität
gefühllos und gelähmt. Der Patient kann ein Kribbeln
empfinden (\E{\enquote{Ameisenlaufen}}).\ZFootnote{Im
englischen werden die Symptome als die
\Speech{\enquote{five P's of limb ischemia}}
zusammengefasst: {Pain} (schmerzend)
{Pale} (blaß),
{Pulseless}  (pulslos),
{Paresthesia} (gefühllos),
{Paralysis}} (Lähmung).

}{%
\begin{compactitem}
  \item Schmerz
  \item Blässe
  \item Kein Puls
  \item Gefühllosigkeit
  \item Lähmung
\end{compactitem}
}{}{}{}


% M %%%%%%%%%%%%%%%%%%%%%%%%%%%%%%%%%%%%%%%%%%%%%%%%%%% M %
% Arterieller Gefäßverschluss einer Extremität
\ProcedurePrintDefault{MI74041C}

\subsubsection{Peripherer \HeadingKeyword{venöser Gefäßverschluss}: 
Tiefe Beinvenenthrombose}
\label{topic:venenthrombose}
\label{topic:thrombose}
\index{Thrombose}
\index{Venenthrombose!tiefe}
\index{Beinvenenthombose}
\index{tVT}




\TextSlide{{\TxBeschreibung} und Ursachen}
{%
\DefinitionInPara{Eine tiefe \T{Beinvenenthrombose} ist
eine Verstopfung einer tiefen Beinvene durch ein
Blutgerinnsel (Thrombus\footnote{Thrombus: Blutgerinnsel. Im Gegensatz dazu 
\T{Embolus}: Losgelöster Thrombus, welcher woanders ein
Gefäß verstopft und eine \T{Embolie} verursacht.}).}
Ursächlich für die Bildung eines Blutgerinnsels ist meistens
eine verminderte Blutflussgeschwindigkeit bzw. Rückstau,
eine Verminderung des flüssigen Anteils des Blutes oder eine gesteigerte bzw.
krankhafte Neigung zur Bildung von Blutgerinnsel. Die \EE{venöse
Insuffizienz} ist ein weit verbreitetes chronisches Leiden.
 Es kommt dabei zu einem Rückfluss von Blut und Aussackungen
 von Venen, welche im fortgeschrittenen Stadium als hässliche
\E{Krampfadern} deutlich sichtbar sind. Schuld sind oft
insuffiziente Venenklappen, stehende bzw. sitzende
Tätigkeiten und fehlende Bewegung.

}{
\begin{itemize}
  \item verminderte Blutflussgeschwindigkeit oder Rückstau
  \item Flüssige Blutbestandteile
  $\downarrow$ 
  \item Gesteigerte bzw. krankhafte Neigung
\end{itemize}
}{}{}{}



\TextSlide
{\TxABCDE}
{~
\begin{ABCDEText}
  \item[\ABCDE{1}] Es finden sich evtl.
  Hinweise auf Bettlägrigkeit oder stattgehabte Reisen
  (Koffer, \ldots). 
  \item[\ABCDE{2}] Evtl. fällt eine
  \EE{Schwangerschaft} auf.
  \item[\ABCDE{4}] \EE{Schmerzen} in der
  betroffenen Extremität.
  \item[\ABCDE{\ldots}] Die betroffene Extremität ist
  schmerzhaft, oft überwärmt, und eher
  rötlich-bläulich-violett (livide) verfärbt. Manchmal sind
  diese Zeichen nicht sehr ausgeprägt.
  In der Regel ist nur eine Extremität betroffen, d.h. die
Symptome sind nur auf \E{einer Seite} zu finden.
\end{ABCDEText}
}
{
\begin{ABCDESlide}
  \item[\ABCDE{1}] Bettlägrigkeit, stattgehabte Reisen
  \item[\ABCDE{2}] Schwangerschaft
  \item[\ABCDE{4}] Schmerzen
  \item[\ABCDE{\ldots}] 
  \begin{compactitem}
  \item Dunkel-bläuliche (livide) Verfärbung. 
  \item Überwärmung
  \item Symptome i.\,d.\,R. \E{einseitig}
  \end{compactitem}
\end{ABCDESlide}
}
{}
{}
{}

\TextSlide
{{\TxSAMPLER}}
{
\begin{SAMPLERText}
  \item[\SAMPLER{S}] Siehe ABCDE.
  \item[\SAMPLER{M}] Eventuell wurden gerinnungshemmende
  Medikamente (\enquote*{Blutverdünnung}:
  Thrombosespritzen, Marcoumar{\texttrademark}, \ldots)
  verschrieben. Eventuell wurden diese nicht oder fehlerhaft
  genommen.
  \item[\SAMPLER{P}] Eventuell sind Krampfadern
  \Z{\I[Varizen!Thrombose]{Varizen})} bekannt. Eventuell
  sind Gerinnungsstörungen bekannt \Z{(z.\,B. \I{Thrombophilie},
  \I{Faktor-V-Leiden-Mutation}, \ldots)}. Weiters sind
  Tumorerkrankungen oder eine Schwangerschaft zu nennen.
  \item[\SAMPLER{L}] Eventuell ist ein vorangehender
  Spitalsaufenthalt, welchem eine Immobilsation folgte,
  erhebbar (Operationen, \ldots).
  \item[\SAMPLER{E}] Verdächtig sind ausgedehnte stehende
  oder sitzende Tätigkeiten (Reisen, \ldots) und
  Bettlägrigkeit, siehe {\SAMPLER{R}}.
  \item[\SAMPLER{R}]
  Problematisch sind insuffiziente Venenklappen, stehende bzw. sitzende
  Tätigkeiten und fehlende Bewegung.
  Eine länger dauernde \EE{Immobilisation} oder
  \EE{Bettlägrigkeit} sind demnach wesentliche
  Risikofaktoren,
  darunter fallen auch lange Reisen ohne Bewegung: So werden
  Thrombosen aufgrund von langen Flugreisen auch scherzhaft als \T{Economy class
  Syndrome} bezeichnet.
  Erschwerend wirkt sich auch eine \EE{unzureichende
  Flüssigkeitsaufnahme} aus, da dann der flüssige Anteil des
  Blutes vermindert sein kann und das Blut folglich
  dickflüssiger wird.
  Ein besonderer Risikofaktor ist der
  \II[Rauchen!Thrombose]{Zigarettenkonsum}.
  Die Einnahme von \EE{hormonellen Verhütungsmitteln}
  (\I[Pille!Thrombose]{Pille}) erhöht zusammen mit
  Zigarettenkonsum das Thromboserisiko drastisch.
  Ebenso erhöht eine bereits \EE{stattgehabte
  Thrombose} das Risiko, wieder eine Thrombose zu bekommen.
  Im Rahmen einer
  \II[Schwangerschaft!Thrombose]{Schwangerschaft} oder von
  Erkrankungen, z.\,B.
  Tumorerkankungen, kann das Gerinnungsystem gestört sein, wodurch es zu Thrombosen
  kommen kann. Ausserdem können Erkrankungen des
  Gerinnungssystems (\SAMPLER{P})  erhebliche
  Risikofaktoren darstellen.
\end{SAMPLERText}
}
{
\begin{SAMPLERSlide}
  \item[\SAMPLER{S}] Siehe {\TxABCDE}
  \item[\SAMPLER{M}] Gerinnungshemmende
  Medikamente -- Einnahmefehler?
  \item[\SAMPLER{P}] Siehe {\SAMPLER{R}}
  \item[\SAMPLER{L}] Vorangehender
  Spitalsaufenthalt
  \item[\SAMPLER{E}] Siehe {\SAMPLER{R}}
  \item[\SAMPLER{R}] Insuffiziente Venenklappen, stehende bzw. sitzende
  Tätigkeiten, Immobilisation, unzureichende
  Flüssigkeitsaufnahme, Rauchen, hormonellen Verhütungsmittel, 
  Thrombose in der Anamnese, Schwangerschaft, 
  Tumorerkankungen, Erkrankungen des
  Gerinnungssystems 
\end{SAMPLERSlide}
}
{}
{}
{}

\TextSlide{\TxGefahren}
{
Es kann zu einer \E{Loslösung von Teilen des Thrombus}
kommen, dieser losgelöste Teil wird
\T{Embolus}\index{Embolus} genannt. Mit dem Blutstrom
gelangt der Embolus über die Hohlvenen in das rechte Herz.
Bis hierhin gibt es noch keine Probleme, da die Blutgefäße
an Dicke zunehmen. Vom rechten Herzen gelangt der
Embolus in den Lungenkreislauf (Lungenarterien). Hier
verzweigen sich die Geäße wieder und der Durchmesser
wird kleiner. Dies führt dazu, dass der Embolus irgendwann
zu groß ist und stecken bleibt. Dadurch wird nun ein 
Lungengefäß verstopft und ein Teil der Lunge nimmt nicht
mehr am Gasaustausch teil: Es entsteht eine
\EE{Lungenembolie}
(\prettyref{topic:lungenembolie}).
}{
\begin{itemize}
  \item Lösung des Thrombus →
  \EE{Lungenembolie} (\prettyref{topic:lungenembolie})
\end{itemize}
}{}{}{}




\begin{PictureSeriesWide}{}{Bilderserie: Thromboserisiken}
\PictureSeriesRowWide{3}
{./images/economyclass.jpg}
{Die Economyclass. Sorgsam geschlichtet verbringen Menschen
hier Stunden damit, Thrombosen zu basteln.
\SourcePicture{Gabriel}}
{./images/lovenox1.jpg}
{Throm\-bose\-prophy\-laxe. \Z{Niedermolekulares Heparin
(hier \enquote{Lovenox\texttrademark}) verhindert Thrombosen die z.\,B. durch lange
Immobilisation (Reisen, Bett\-lägrig\-keit, Gips, \ldots)
entstehen können. Die Substanz wird unter die Haut (\enquote{subkutan}) gespritzt.
\SourcePicture{Gabriel}}}
{./images/hirtler-aass/Thrombus-Embolie}
{Eine Venenthrombose kann eine Embolie auslösen. 
\SourcePicture{Hirtler}}
{empty}
{}
\end{PictureSeriesWide}


% M %%%%%%%%%%%%%%%%%%%%%%%%%%%%%%%%%%%%%%%%%%%%%%%%%%% M %
% Venöser Gefäßverschluss einer Extremität
\ProcedurePrintDefault{MI74042C}

\citep{Renz-Polster2006,SiewertChirurgie8,ScholzNotfallmedizin2,IntroClinEmergMed4,Longmore2007,Fauci2008}

% \paragraph{Weiters}
\nocite{Eisenburger2008,
Havel2007,
AMLS3,
Erc2005DeSec5,
Erc2005Sec5,
Erc2005DeSec6,
Erc2005DeSec3,
Erc2005Sec3,
Juurlink2005,
Kalla2006,
Vlcek2008,
Wells2000,
BoehmerTaschRett6,
LPN3,
Boecker2,
Helfen2008,
LippertAnatomie7,
Patzelt2005,
Sachsenweger2003,
RedelsteinerHRNS1,
ForMed2,
StriebelAn7,
Zeiler2001,
KlinkePhysio3,
DRAn3,
ForthPharma8,
LuellmannPharma16,
ThierbachInterhospitaltransfer1,
AnCompact3,
RueckerBildatlas1}

%\putbib
% % Rea während Transport
% \nocite{Eisenburger2008,Havel2007}
% % General
% \nocite{SiewertChirurgie8}
% \nocite{AMLS3}
% \nocite{Erc2005DeSec5}
% \nocite{Erc2005Sec5}
% \nocite{Erc2005DeSec6}
% \nocite{Erc2005DeSec3}
% \nocite{Erc2005Sec3}
% \nocite{Juurlink2005}
% \nocite{Kalla2006}
% 
% \nocite{Vlcek2008}
% \nocite{Wells2000}
% \nocite{BoehmerTaschRett6}
% \nocite{Fauci2008}
% \nocite{Lissauer2007}
% 
% \nocite{LPN3,Boecker2}
% \nocite{Helfen2008}
% \nocite{LippertAnatomie7}
% \nocite{Nilsson1997}
% \nocite{Patzelt2005}
% \nocite{Sachsenweger2003}
% \nocite{RedelsteinerHRNS1}
% \nocite{ForMed2}
% \nocite{StriebelAn7}
% \nocite{Zeiler2001}
% \nocite{KlinkePhysio3}
% \nocite{DRAn3}
% \nocite{ForthPharma8}
% \nocite{LuellmannPharma16}
% \nocite{ThierbachInterhospitaltransfer1}
% \nocite{AnCompact3}
% \nocite{RueckerBildatlas1}


% \clearpage
\ChapterBibliography 

%\bibliographystyle{unsrtdin2}{\scriptsize \bibliography{Literatur-AAS}}

\ChapterEnd